\documentclass[11pt,a4paper]{report}

\usepackage[T1]{fontenc}
\usepackage[utf8]{inputenc}
\usepackage{amsmath,amssymb,amsthm}
\usepackage{booktabs}
\usepackage{longtable}
\usepackage{array}
\usepackage{microtype}
\usepackage[margin=3.2cm]{geometry}
\usepackage{xcolor}
\usepackage[colorlinks,linkcolor=blue!50!black,citecolor=blue!50!black,urlcolor=blue!50!black]{hyperref}
\usepackage[normalem]{ulem}

\theoremstyle{plain}
\newtheorem{theorem}{Theorem}[chapter]
\newtheorem{proposition}[theorem]{Proposition}
\newtheorem{lemma}[theorem]{Lemma}
\newtheorem{corollary}[theorem]{Corollary}
\theoremstyle{definition}
\newtheorem{definition}[theorem]{Definition}
\newtheorem{example}[theorem]{Example}
\theoremstyle{remark}
\newtheorem{remark}[theorem]{Remark}


\newcommand{\E}{\operatorname{E}}
\newcommand{\Umbral}{\mathsf{Umbral}}
\newcommand{\sK}{\mathbb{K}}
\newcommand{\sU}{\mathbb{U}}
\newcommand{\sA}{\mathbb{A}}
\newcommand{\NN}{\mathbb{N}}
\newcommand{\ZZ}{\mathbb{Z}}
\newcommand{\QQ}{\mathbb{Q}}
\newcommand{\RR}{\mathbb{R}}
\newcommand{\T}{\mathsf{T}}
\newcommand{\Nat}{\mathsf{N}}
\newcommand{\W}{\mathcal{W}}
\newcommand{\mk}{\mathrm{mk}}
\newcommand{\den}[2]{[\![#1]\!]_{#2}}
\newcommand{\rel}[2]{\langle\!\langle #1 \rangle\!\rangle_{#2}}
\newcommand{\by}[1]{\mathrel{\overset{\textup{(#1)}}{=}}}

\setlength{\emergencystretch}{3em}

\title{Free Theorems for Umbral Calculus}
\author{Sonat S\"uer}
\date{\today}

\begin{document}
\maketitle

\begin{abstract}
These notes are about an implementation of umbral calculus in Haskell as a library and 
the free theorems about umbral operations. It turns out that many standard assumptions which are proved
case by case in the literature can be deduced as free theorems in the sense of \cite{Wadler1989} in a unified way. The results are
formalized in Lean with the help of Claude Opus~5.5. Claude generated most of the proofs though the
main definitions are human made. Parts of the Haskell code (most notably the test suite, the showcase 
and the general purpose libraries) were initially generated by Claude but later revised.
\end{abstract}

\tableofcontents

\chapter{Basic Notions of Umbral Calculus}
\label{ch:intro}

The first chapter is a fast but hopefully intuitive presentation of basic 
definitions of umbral calculus.  

\section{The Bernoulli Recurrence by Magic}
\label{sec:blissard}

Umbral calculus is a branch of algebraic combinatorics. Historically it stems from
the observation that certain polynomial identities resemble one another for no
apparent reason. Using these resemblances, one can produce correct results with unjustified proofs.
Here is an example to give you an idea: One can pretend that the subscript of the Bernoulli
number $B_n$ is an exponent, manipulate $B^n$ as if $B$ were a number,
and turn the exponents back to subscripts to obtain results about Bernoulli numbers. This \sout{utterly insane}
symbolic method goes back to Blissard in the nineteenth century \cite{Blissard1861}. 

Let us see a concrete example. Recall that the Bernoulli numbers
$B_n$ are defined by
\[
  \frac{t}{e^t - 1} = \sum_{n \geq 0} B_n \frac{t^n}{n!} .
\]
It is a classical and easy fact that they satisfy the recurrence relation
\begin{equation}
\label{eq:bernoulli-rec}
  \sum_{i=0}^{n-1} \binom{n}{i} B_i = 0 \qquad (n \geq 2) .
\end{equation}
The magical proof of this fact writes $B^n$ for $B_n$ in the defining equation. The defining
series then becomes $e^{Bt}$ because 
\[
\frac{t}{e^t - 1} =
\sum_{n \geq 0} B_n \frac{t^n}{n!} = 
\sum_{n \geq 0} B^n \frac{t^n}{n!} =
\sum_{n \geq 0} \frac{(Bt)^n}{n!} =
e^{Bt}.
\]
This gives us
\[
  t = e^{Bt}(e^t - 1) = e^{(B+1)t} - e^{Bt} .
\]
The coefficient of $t^n/n!$ on the left-hand side is $1$ for $n = 1$ and $0$ otherwise, 
which means a lot of things simplify on the right-hand side. More precisely, $(B+1)^n = B^n$ 
holds for $n \geq 2$. So, for $n \geq 2$, we have
\[
B^n = (B + 1)^n = \sum_{i=0}^n \binom{n}{i} B^i = B^n + \sum_{i=0}^{n - 1} \binom{n}{i} B^i .
\]
Finally, canceling out $B^n$ on both sides and writing $B_n$ for $B^n$ gives us 
exactly~\eqref{eq:bernoulli-rec}. Now let us make sense of this.


\section{Formal Definition of Umbral Calculus}
\label{sec:umbrae}

Rota and Taylor made calculations like the one in the previous section rigorous in a very elegant and 
deceptively simple way in \cite{RotaTaylor1994}. For a more contemporary account see \cite{Gessel2003}.
They define an umbral calculus by the following data:

\begin{definition}
\label{def:umbral}
An \emph{umbral calculus} consists of a commutative ring $k$, a set $\mathcal{A}$
of symbols called \emph{umbrae}, and a $k$-linear map
$\E : k[\mathcal{A}] \to k$, the \emph{evaluation}, such that
\begin{enumerate}
\item[(U1)] $\E[1] = 1$;
\item[(U2)] (uncorrelation) for distinct umbrae $\alpha, \beta, \dots, \gamma$ and
  exponents $i, j, \dots, l$,
  \[
    \E[\alpha^{i}\beta^{j}\cdots\gamma^{l}]
      = \E[\alpha^{i}]\,\E[\beta^{j}]\cdots\E[\gamma^{l}] .
  \]
\end{enumerate}
The \emph{moments} of an umbra $\alpha$ are $a_n = \E[\alpha^n]$.
Two umbrae are \emph{similar}, written $\alpha \equiv \beta$, if they have the
same moments. Rota and Taylor call them \emph{exchangeable}; we follow
Di Nardo and Senato \cite{DiNardoSenato2001}.\footnote{Rota and Taylor also require an augmentation umbra, all of whose moments of
positive order vanish, but we will not need it in these notes.}
\end{definition}

This is an axiomatic definition and it does not mention any models. In these notes we will talk about two models.
One is a probabilistic model in which an umbra is a random variable, the evaluation is expectation,
and distinct umbrae are independent. The other model is based on moments. We will cover both of them in~\ref{sec:moment}.

Note that uncorrelation is not multiplicativity, that is, evaluation is not a ring homomorphism. For an umbra $\alpha$
with moments $a_n$ we have $\E[\alpha^2] = a_2$ and $\E[\alpha]^2 = a_1^2$. Under the probabilistic interpretation this means
that one can talk about the variance of an umbra defined by $\E[\alpha^2] - \E[\alpha]^2$.

We need one more definition before giving a proper proof for~\eqref{eq:bernoulli-rec}.

\begin{definition}
\label{def:gf}
The \emph{coefficientwise extension} of the evaluation is the $k$-linear map
\[
  \E : k[\mathcal{A}][[t]] \to k[[t]] ,
  \qquad
  \E\Bigl[\sum_{n \geq 0} p_n t^n\Bigr] = \sum_{n \geq 0} \E[p_n]\, t^n ,
\]
where $t$ is a formal variable. For convenience we denote this function also by $\E$.
\end{definition}

Let us assume that $k$ contains $\QQ$. Then for $p \in k[\mathcal{A}]$ we can talk about the series $e^{pt}$,
which is given by $\sum_{n \geq 0} \frac{p^n}{n!}t^n$. Since $k[\mathcal{A}]$ is a commutative $\QQ$-algebra, 
one can prove that the identity $e^{(p+q)t} = e^{pt}e^{qt}$ holds for all $p$ and $q$ in $k[\mathcal{A}]$.

With these new definitions and tools now we can justify Blissard's trick. 

\section{The Bernoulli Recurrence by Sufficiently Advanced Technology}
\label{sec:bernoulli}
Assume that $k$ contains $\QQ$.

\begin{definition}
\label{def:bernoulli}
A \emph{Bernoulli umbra} is an umbra $\iota$ satisfying $\E[\iota^n] = B_n$ for all $n$.
\end{definition}

Note that this is equivalent to the identity 
\[
  \E[e^{\iota t}] = \frac{t}{e^t - 1} .
\]

\begin{proposition}
\label{prop:bernoulli-rec}
Let $\iota$ be a Bernoulli umbra. Then $\E[(\iota+1)^n] = \E[\iota^n]$ for
$n \geq 2$. Equivalently, the Bernoulli numbers satisfy~\eqref{eq:bernoulli-rec}.
\end{proposition}

\begin{proof}
By the above remark we have
\begin{align*}
  \E[e^{(\iota+1)t}] - \E[e^{\iota t}]
    &= \E[e^{t} e^{\iota t}] - \E[e^{\iota t}] \\
    &= e^{t}\,\E[e^{\iota t}] - \E[e^{\iota t}] \\
    &= (e^{t} - 1)\,\frac{t}{e^{t} - 1} \\
    &= t .
\end{align*}
Comparing the coefficients of $t^n/n!$ gives $\E[(\iota+1)^n] = \E[\iota^n]$ for
$n \geq 2$. On the other hand, the binomial theorem in $k[\mathcal{A}]$ yields
\[
  \E[(\iota+1)^n]
    = \E\Bigl[\,\sum_{i=0}^{n} \binom{n}{i} \iota^i\Bigr]
    = \sum_{i=0}^{n} \binom{n}{i} \E[\iota^i]
    = \sum_{i=0}^{n} \binom{n}{i} B_i .
\]
Subtracting $\E[\iota^n] = B_n$ from both sides gives~\eqref{eq:bernoulli-rec}.
\end{proof}

This is lovely, but it doesn't showcase all the main ideas of umbral calculus. It only uses the
fact that $\E$ is linear. For some identities we need more than linearity. Here is an example.

\section{Products of Bernoulli Numbers and the Role of Similarity}
\label{sec:euler}

Another classical identity satisfied by Bernoulli numbers is Euler's identity:
\begin{equation}
\label{eq:euler}
  \sum_{i=0}^{n} \binom{n}{i} B_i B_{n-i} = -n B_{n-1} - (n-1) B_n
  \qquad (n \geq 1) .
\end{equation}
Write $f(t) = t/(e^t - 1)$. Differentiating gives
\[
f'(t) = 
\frac{e^t -1 - te^t}{(e^t - 1)^2} = 
\frac{1}{t} \left( \frac{t(e^t - 1)}{(e^t - 1)^2} - \frac{e^tt^2}{(e^t - 1)^2} \right) = 
\frac{1}{t}(f(t) - e^tf(t)^2)
\]
which is equivalent to $t f'(t) = f(t) - f(t)^2 e^t$. To get rid of the last $e^t$, we observe that $f(t)^2(e^t - 1) = t f(t)$, which yields
\begin{equation}
\label{eq:fsquare}
  f(t)^2 = f(t) - t f(t) - t f'(t) .
\end{equation}
The coefficients of $t^n/n!$ in $f$, $tf$ and $tf'$ are $B_n$, $n B_{n-1}$
and $n B_n$, respectively. So the identity~\eqref{eq:euler} says that the coefficient of
$t^n/n!$ in $f^2$ is the sum on its left. So far so good. We only used simple algebra on the generating function. 

If we try to use Blissard's method to proceed, we get into trouble because since $f(t)^2 = e^{Bt}e^{Bt}$
by the magical method, we should have $f(t)^2 = e^{2Bt}$, which has the wrong coefficients,
namely $2^n B_n$. Apparently we need more magic. So we \emph{invent} another symbol $B'$, which behaves 
like $B$ but is not $B$. More precisely, $B^i B'^j$ stands for $B_i B_j$. Note that $'$ disappears 
after turning the exponents to indices. Then $f(t)^2 = e^{(B+B')t}$ and the solution follows. Let us 
write this out in detail marking the assumptions we need to make. We already have two
isolated assumptions, namely linearity (Lin) and uncorrelation (Unc) coming from the definition. 

\begin{proposition}
\label{prop:euler}
Let $\iota$ be a Bernoulli umbra, and let $\iota' \equiv \iota$ be an umbra
distinct from $\iota$. Then
\[
  \E\bigl[(\iota+\iota')^n\bigr]
    = \sum_{i=0}^{n} \binom{n}{i} B_i B_{n-i}
    = -n B_{n-1} - (n-1) B_n
  \qquad (n \geq 1) .
\]
\end{proposition}

\begin{proof}
We have
\[
  f(t)^2
    = \E[e^{\iota t}]\,\E[e^{\iota t}]
    \by{?} \E[e^{\iota t}]\,\E[e^{\iota' t}]
    \by{Unc} \E[e^{\iota t} e^{\iota' t}]
    = \E[e^{(\iota+\iota')t}] .
\]
Comparing the coefficients of $t^n/n!$ with~\eqref{eq:fsquare} gives
$\E[(\iota+\iota')^n] = -n B_{n-1} - (n-1) B_n$. On the other hand,
\[
  \E\bigl[(\iota+\iota')^n\bigr]
    \by{Lin} \sum_{i=0}^{n} \binom{n}{i} \E[\iota^i \iota'^{\,n-i}]
    \by{Unc} \sum_{i=0}^{n} \binom{n}{i} \E[\iota^i]\,\E[\iota'^{\,n-i}]
    = \sum_{i=0}^{n} \binom{n}{i} B_i B_{n-i} .
\]
\end{proof}

In the computations above, the use of (Lin) and (Unc) is clear. But (?) is new. 
In the step marked by (?), what we do is to substitute one umbra with a similar one.
This principle is usually stated in the following form: 

\begin{quote}
\textbf{Substitution principle (Sub).} If $\alpha \equiv \beta$, then $\alpha$ may be
replaced by $\beta$ in any umbral expression without changing its value,
provided $\beta$ does not already occur in it.
\end{quote}

But is (Sub) a sound principle? In this particular case the answer is yes and it is
easy to prove. The general case is a little ambiguous because the principle says \emph{any umbral expression}
and we do not have a definition of an umbral expression. These notes clarify this issue.
They define a domain-specific language for umbral calculus in which (Sub) is
captured as a free theorem.


\section{Two Models of Umbral Calculus}
\label{sec:moment}

The first model is so straightforward that it almost has no content. The reason I am 
recording here is that the Haskell implementations uses it as one of its backends.

\begin{definition}
\label{def:moment}
Let $k$ be a commutative ring and $\mathcal{A}$ a set, whose elements are called
\emph{names}. An \emph{assignment} is a function
$\rho : \mathcal{A} \to k^{\NN}$. For a name $\alpha$ and $n \geq 1$, the term
$\rho(\alpha)_n$ of the sequence $\rho(\alpha)$ is the \emph{$n$-th moment} of
$\alpha$ under $\rho$. The \emph{moment model} of $\rho$ consists of $k$,
$\mathcal{A}$ and the $k$-linear map $\E_\rho : k[\mathcal{A}] \to k$ defined on
monomials by
\[
  \E_\rho\Bigl[\,\prod_{\alpha \in \mathcal{A}} \alpha^{e_\alpha}\Bigr]
    = \prod_{\alpha : e_\alpha > 0} \rho(\alpha)_{e_\alpha} ,
\]
where the exponents $e_\alpha$ are zero for all but finitely many $\alpha$, and the
empty product is $1$.
\end{definition}

The $0$-th moment is $\E_\rho[1] = 1$ whatever $\rho$ is.
The term $\rho(\alpha)_0$ is never used so one might as well use $k^{\NN^+}$ as the target of assignments. 


\begin{proposition}
\label{prop:moment}
For every assignment $\rho$, the moment model of $\rho$ is an umbral calculus, in
which $\E_\rho[\alpha^n] = \rho(\alpha)_n$ for every name $\alpha$ and every
$n \geq 1$. Conversely, every umbral calculus $(k, \mathcal{A}, \E)$ is the moment
model of the assignment $\rho(\alpha)_n = \E[\alpha^n]$.
\end{proposition}

\begin{proof}
The monomial $1$ has all exponents zero, so $\E_\rho[1]$ is the empty product,
which gives (U1). For distinct names $\alpha, \beta, \dots, \gamma$, the right
side of the definition at $\alpha^i\beta^j\cdots\gamma^l$ is
$\rho(\alpha)_i\,\rho(\beta)_j\cdots\rho(\gamma)_l$, and each factor is the value
of $\E_\rho$ at the corresponding single power. This gives (U2) and the formula
for the moments. Conversely, by (U2) the evaluation $\E$ of a monomial is the
product of the moments of its umbrae, which is the value of $\E_\rho$ at it. Both
maps are $k$-linear and agree on monomials, so $\E = \E_\rho$.
\end{proof}

Interpreting umbrae as random variables gives a different model of umbral calculus.

\begin{definition}
\label{def:probmodel}
Let $(\Omega, \mu)$ be a probability space, $\mathcal{A}$ a set of names, and
$X = (X_\alpha)_{\alpha \in \mathcal{A}}$ a family of real functions on $\Omega$.
For an assignment $\rho$ over $\RR$, the family $X$ is a \emph{model} of $\rho$ if
it is independent, each $X_\alpha$ is measurable with every power integrable, and
$\int X_\alpha^n\, d\mu = \rho(\alpha)_n$ for every name $\alpha$ and every
$n \geq 1$. For $p \in \RR[\mathcal{A}]$, let $p(X)$ be the function
$\omega \mapsto p\bigl(X(\omega)\bigr)$ obtained by substituting $X_\alpha$ for
$\alpha$. The \emph{probabilistic model} of $X$ consists of $\RR$, $\mathcal{A}$
and the map $\E_X : \RR[\mathcal{A}] \to \RR$ defined by
\[
  \E_X[p] = \int p(X)\, d\mu .
\]
\end{definition}

\begin{proposition}
\label{prop:probmodel}
Let $X$ be a model of an assignment $\rho$. Then the probabilistic model of $X$ is
the moment model of $\rho$, that is, $\E_X = \E_\rho$. In particular, it is an
umbral calculus.
\end{proposition}

\begin{proof}
Consider a monomial $\prod_\alpha \alpha^{e_\alpha}$. By independence, the
product $\prod_\alpha X_\alpha^{e_\alpha}$ of integrable functions is integrable,
and its integral is the product of the integrals
$\int X_\alpha^{e_\alpha}\, d\mu = \rho(\alpha)_{e_\alpha}$ of its factors with
$e_\alpha > 0$. Hence $\E_X$ and $\E_\rho$ agree on monomials. Every $p(X)$ is a
finite linear combination of such products, so it is integrable, and both maps
are $\RR$-linear. Thus $\E_X = \E_\rho$, which is an umbral calculus by
Proposition~\ref{prop:moment}.
\end{proof}

\chapter{Umbral Calculus in a Fragment of System F}
\label{ch:restricted}

This chapter defines a fragment of System~F and gives a rather straightforward implementation 
of umbral calculus in it. Working in a small fragment allows us to define a set theoretic 
semantics in which we can prove a version Reynold's abstraction theorem. It is well known 
that there is no such semantics for full Sytem~F, \cite{Reynolds1984}. Another advantage of 
a small language is that proofs are simple.

\section{A Restricted System F}
\label{sec:calculus}

By restricted I mean two things: Quantifiers range over quantifier-free types, so the fragment is
\emph{predicative}. A type abstraction may only be formed over a term with no free
term variables, so quantifiers stand in front, and the fragment is
\emph{prenex}. Due to lack of imagination, I will call it the \emph{restricted language}. 
I should also mention that nothing is original here. This section exists essentially to fix notation.

\begin{definition}[Types]
\label{def:types}
Fix a set of \emph{type variables} $X, Y, \dots$. The \emph{types} are
generated by
\[
  \tau \;::=\; X \mid \tau \to \tau \mid \forall X.\, \tau .
\]
\end{definition}

\begin{definition}[Terms]
\label{def:terms}
A \emph{context} $\Gamma$ is a
finite list of distinct term variables, each with a type. A \emph{judgement}
$\Gamma \vdash M : \tau$ consists of a context $\Gamma$, a term $M$ and a
type $\tau$. The derivable judgements are those generated by the rules
\begin{gather*}
  \frac{x : \tau \in \Gamma}{\Gamma \vdash x : \tau}
  \qquad
  \frac{\Gamma, x : \sigma \vdash M : \tau}
       {\Gamma \vdash \lambda x.\, M : \sigma \to \tau}
  \\[1ex]
  \frac{\Gamma \vdash M : \sigma \to \tau \quad \Gamma \vdash N : \sigma}
       {\Gamma \vdash M\,N : \tau}
  \qquad
  \frac{\varnothing \vdash M : \tau}
       {\varnothing \vdash \Lambda X.\, M : \forall X.\, \tau}
\end{gather*}
A \emph{closed term} of type $\tau$ is a term $M$ with
$\varnothing \vdash M : \tau$.
\end{definition}

Now the semantics.

\begin{definition}[Semantics]
\label{def:semantics}
Fix a set $\mathcal{B}$ of \emph{base sorts}. The \emph{simple types}
are generated by
\[
  \sigma \;::=\; b \mid \sigma \to \sigma \qquad (b \in \mathcal{B}).
\]
An \emph{interpretation} is a function $I$ that assigns a set $I(b)$ to every base
sort $b$. The \emph{denotation} of a simple type in $I$ is the set
$\den{\sigma}{I}$ defined by
\[
  \den{b}{I} = I(b) ,
  \qquad
  \den{\sigma \to \sigma'}{I} = \den{\sigma'}{I}^{\,\den{\sigma}{I}} .
\]
A \emph{type environment} is a function $\eta$ from type variables to simple
types. For a type environment $\eta$, a type variable $X$ and a simple type
$\sigma$, write $\eta[X \mapsto \sigma]$ for the type environment that sends $X$ to
$\sigma$ and agrees with $\eta$ at every other type variable. The denotation of a type in $I$ and $\eta$ is the set
$\den{\tau}{I,\eta}$ defined by the same clause at $\to$, together with
\[
  \den{X}{I,\eta} = \den{\eta(X)}{I} ,
  \qquad
  \den{\forall X.\, \tau}{I,\eta}
    = \prod_{\sigma \text{ simple}} \den{\tau}{I,\eta[X \mapsto \sigma]} .
\]
The denotation of a term $\Gamma \vdash M : \tau$ is the function
$\den{M}{I,\eta}$ from $\den{\Gamma}{I,\eta}$, the product of the denotations of
the types in $\Gamma$, to $\den{\tau}{I,\eta}$, defined by induction on the
derivation in the standard way. In particular $\den{\Lambda X.\, M}{I,\eta}$ is
the family
$\bigl(\den{M}{I,\eta[X \mapsto \sigma]}\bigr)_{\sigma}$.
\end{definition}

The product over simple types is a set because the simple types form a set. This
is where predicativity is used.

\begin{definition}[Logical Relation]
\label{def:logrel}
Let $I$ and $J$ be interpretations. A \emph{relation
environment} is a triple $\xi = (\eta, \eta', \varrho)$, where $\eta$ and $\eta'$
are type environments and $\varrho$ is a function on type variables with
$\varrho(X) \subseteq \den{\eta(X)}{I} \times \den{\eta'(X)}{J}$ for every $X$.
For simple types $\sigma$, $\sigma'$ and a relation
$Q \subseteq \den{\sigma}{I} \times \den{\sigma'}{J}$, write
$\xi[X \mapsto (\sigma, \sigma', Q)]$ for
$(\eta[X \mapsto \sigma], \eta'[X \mapsto \sigma'], \varrho[X \mapsto Q])$.
The \emph{logical relation} at a type $\tau$ is the relation
\[
  \tau(\xi) \subseteq \den{\tau}{I,\eta} \times \den{\tau}{J,\eta'}
\]
defined by induction on $\tau$ as follows.
\begin{itemize}
\item $X(\xi) = \varrho(X)$.
\item $f \mathrel{(\tau \to \tau')(\xi)} g$ if and only if
  $f(x) \mathrel{\tau'(\xi)} g(y)$ for all $x \mathrel{\tau(\xi)} y$.
\item $f \mathrel{(\forall X.\, \tau)(\xi)} g$ if and only if
  $f(\sigma) \mathrel{\tau(\xi[X \mapsto (\sigma, \sigma', Q)])} g(\sigma')$
  for all simple types $\sigma$, $\sigma'$ and every relation
  $Q \subseteq \den{\sigma}{I} \times \den{\sigma'}{J}$.
\end{itemize}
If $\tau$ has no free type variables, then $\tau(\xi)$ does not depend on
$\xi$, and it is written $\tau(\varnothing)$.
\end{definition}

\begin{theorem}[Abstraction Theorem]
\label{thm:abstraction-precise}
Let $I$ and $J$ be interpretations. Then the denotations in $I$ and in $J$ of
every closed term of type $\tau$ are related by $\tau(\varnothing)$.
\end{theorem}

\begin{proof}
By induction on the derivation of $\Gamma \vdash M : \tau$, for related
environments, as in Reynolds \cite{Reynolds1983}. The case of $\Lambda X.\, M$
uses that the relation at $\forall X.\, \tau$ quantifies over every relation $Q$.
\end{proof}

\section{Encoding the Umbral Calculus}
\label{sec:encoding}

In order to do umbral calculations, one needs to pass the umbral context, that is the setup
described in~\ref{def:umbral}, to the function doing the actual work. This 
can be done in many ways. I opted for a curried dictionary passing style to keep the language
small. Also, I did not mention polynomials specifically and worked with rings. So, the data 
needed to perform basic umbral calculations is a type $\Umbral(K, U, A)$ defined as below.
\begin{definition}
\label{def:class}
Let $K$, $U$ and $A$ be type variables. The \emph{methods} of the umbral calculus
are
\begin{gather*}
  0_K, 1_K : K ,
  \qquad
  {+_K}, {\cdot_K} : K \to K \to K ,
  \qquad
  {-_K} : K \to K ,
  \\
  0_A, 1_A : A ,
  \qquad
  {+_A}, {\cdot_A} : A \to A \to A ,
  \qquad
  {-_A} : A \to A ,
  \\
  \mathrm{const} : K \to A ,
  \qquad
  \mathrm{inj} : U \to A ,
  \qquad
  \E : A \to K .
\end{gather*}
Here $K$ stands for the \emph{scalars}, $U$ for the \emph{umbrae} and $A$ for
the \emph{umbral polynomials}. For a type $\tau$, the type
$\Umbral(K, U, A) \to \tau$ is the curried function type that takes the methods
in this order,
\begin{multline*}
  K \to K \to (K \to K \to K) \to (K \to K \to K) \to (K \to K) \\
  \to A \to A \to (A \to A \to A) \to (A \to A \to A) \to (A \to A) \\
  \to (K \to A) \to (U \to A) \to (A \to K) \to \tau .
\end{multline*}
For sets $\sK$, $\sU$ and $\sA$, a \emph{dictionary} at $\sK$, $\sU$ and $\sA$ is a tuple
of elements and functions of these types, one for each method.
\end{definition}

An \emph{expression} of the language is a closed term, in the sense of
Definition~\ref{def:terms}, of the polymorphic type
\begin{equation}
\label{eq:type}
  \forall K\,U\,A.\ \Umbral(K, U, A) \to U \to K .
\end{equation}
Here $K$, $U$ and $A$ are type variables. To read a closed term $t$ at a dictionary $d$ with carriers $\sK$,
$\sU$, $\sA$, take three base sorts interpreted as these carriers, instantiate the
three type variables at them, and apply the result to the methods of
$d$. This gives the function $\den{t}{d} : \sU \to \sK$.

More generally, an \emph{expression in $m$ umbrae} is a closed term of the type
\[
  \forall K\,U\,A.\ \Umbral(K, U, A) \to
    \underbrace{U \to \cdots \to U}_{m} \to K ,
\]
and it denotes a function $\den{t}{d} : \sU^m \to \sK$ in the same way. An
expression is an expression in one umbra.

The \emph{moment dictionary} of an assignment $\rho$
(Definition~\ref{def:moment}) has $\sK = k$, $\sU = \mathcal{A}$ the names,
$\sA = k[\mathcal{A}]$, the ring operations of $k$ and of $k[\mathcal{A}]$, the
constant polynomials, the variables, and the evaluation $\E_\rho$. We write
$\den{t}{\rho}$ for the corresponding denotation.

\section{Free Theorems}
\label{sec:free}

After an long setup now we are in a position to prove some free theorems.

\begin{definition}
\label{def:respects}
Let $d_1$ and $d_2$ be dictionaries at $(\sK_1, \sU_1, \sA_1)$ and at
$(\sK_2, \sU_2, \sA_2)$ respectively. Three relations, $R_K$ between $\sK_1$ and $\sK_2$,
$R_U$ between $\sU_1$ and $\sU_2$, and $R_A$ between $\sA_1$ and $\sA_2$, \emph{respect}
$d_1$ and $d_2$ if the two zeros and the two ones are related, each pair of
corresponding operations sends related arguments to related results, and
\begin{gather*}
  x \mathrel{R_K} y \ \Rightarrow\ \mathrm{const}(x) \mathrel{R_A} \mathrm{const}(y),
  \qquad
  u \mathrel{R_U} v \ \Rightarrow\ \mathrm{inj}(u) \mathrel{R_A} \mathrm{inj}(v),
  \\
  p \mathrel{R_A} q \ \Rightarrow\ \E(p) \mathrel{R_K} \E(q) .
\end{gather*}
\end{definition}

\begin{theorem}[The Free Theorem in $m$ Umbrae]
\label{thm:free}
Let $t$ be an expression in $m$ umbrae, and let $R_K$, $R_U$, $R_A$ respect
dictionaries $d_1$ and $d_2$. If $u_i \mathrel{R_U} v_i$ for
$i = 1, \dots, m$, then
$\den{t}{d_1}(u_1, \dots, u_m) \mathrel{R_K} \den{t}{d_2}(v_1, \dots, v_m)$.
\end{theorem}

\begin{proof}
This is Theorem~\ref{thm:abstraction-precise}, read at the dictionary-passing
form of the type of expressions in $m$ umbrae. The three type variables receive
the relations $R_K$, $R_U$ and $R_A$. The dictionary arguments are related by
Definition~\ref{def:respects}, and the umbra arguments by hypothesis.
\end{proof}

This is still very abstract. By instantiating Theorem~\ref{thm:free} at specific relations 
we can obtain more familiar results. Here is an example.

\begin{lemma}[Renaming]
\label{lem:rename}
Let $\kappa$ be an injective map of names such that $\kappa(\alpha)$ has, under
$\rho_2$, the moments that $\alpha$ has under $\rho_1$. Write $\kappa_*$ for the
induced renaming of polynomials. Then
$\E_{\rho_2}[\kappa_*(p)] = \E_{\rho_1}[p]$ for every $p$.
\end{lemma}

\begin{proof}
By linearity it suffices to treat a monomial $\prod_\alpha \alpha^{e_\alpha}$. Its
image is $\prod_\alpha \kappa(\alpha)^{e_\alpha}$, and since $\kappa$ is
injective no two exponents are added together. The evaluation of the image is
therefore the product of the moments
$\rho_2(\kappa(\alpha))_{e_\alpha} = \rho_1(\alpha)_{e_\alpha}$.
\end{proof}

\begin{theorem}[Substitution Principle]
\label{thm:substitution}
Let $\rho$ be an assignment, let $\alpha \equiv \beta$ be two similar umbrae, and
let $\gamma_1, \dots, \gamma_m$ be umbrae different from $\alpha$ and from
$\beta$. Then
\[
  \den{t}{\rho}(\gamma_1, \dots, \gamma_m, \alpha)
    = \den{t}{\rho}(\gamma_1, \dots, \gamma_m, \beta)
\]
for every expression $t$ in $m + 1$ umbrae.
\end{theorem}

\begin{proof}
Let $\theta$ be the transposition exchanging $\alpha$ and $\beta$. Take $R_K$ to
be equality, $R_U$ the graph of $\theta$, and $R_A$ the graph of $\theta_*$. The
conditions for all methods other than $\E$ hold because $\theta_*$ is a ring
homomorphism that fixes the constants and sends the variable $\gamma$ to
$\theta(\gamma)$. The condition for the evaluation is
$\E_\rho[\theta_*(p)] = \E_\rho[p]$, which is Lemma~\ref{lem:rename}: $\theta$ is
injective, and it preserves moments because $\alpha$ and $\beta$ are similar.
Finally $\gamma_i \mathrel{R_U} \gamma_i$, because $\theta$ fixes every umbra
other than $\alpha$ and $\beta$, and $\alpha \mathrel{R_U} \beta$. Theorem~\ref{thm:free}
in $m + 1$ umbrae gives the claim.
\end{proof}

Note that by Proposition~\ref{prop:moment} every
umbral calculus is the moment model of its moments, so the theorem holds in every
umbral calculus.

It is a good point to talk about the uncorrelation condition (U2) in the umbral calculus
Definition~\ref{def:umbral} with regard to free theorems. First of all, it does not follow from
Theorem~\ref{thm:free}. A free theorem gives information about how expressions behave across
related dictionaries. However there is a relation between uncorrealtion and the underlying machinery.

\begin{theorem}[Uncorrelation Constrains Respecting Relations]
\label{thm:uncorrelation}
Let $\rho_1$ and $\rho_2$ be assignments, and let $R_K$, $R_U$, $R_A$ respect the
moment dictionaries of $\rho_1$ and $\rho_2$, with $R_K$ equality on the
scalars. If $R_U$
relates two distinct umbrae $\alpha$ and $\gamma$ to the same umbra $\beta$, then
\[
  \E_{\rho_1}[\alpha^i]\,\E_{\rho_1}[\gamma^j] = \E_{\rho_2}[\beta^{i+j}]
  \qquad (i, j \geq 1) .
\]
In particular, if $\alpha$ and $\gamma$ under $\rho_1$ and $\beta$ under $\rho_2$
have the same moments $a_n$, and $a_n \neq a_1^n$ for some $n$, then $R_U$ does not relate
both $\alpha$ and $\gamma$ to $\beta$.
\end{theorem}

\begin{proof}
By the condition for $\mathrm{inj}$, both $\alpha$ and $\gamma$ are related to
$\beta$ by $R_A$. By the condition for multiplication, applied $i - 1$
times, $\alpha^i$ is related to $\beta^i$, and likewise $\gamma^j$ to $\beta^j$.
Applied once more, it relates $\alpha^i\gamma^j$ to $\beta^{i+j}$. By the condition for the
evaluation, $\E_{\rho_1}[\alpha^i\gamma^j] = \E_{\rho_2}[\beta^{i+j}]$, and by
uncorrelation the left side is $\E_{\rho_1}[\alpha^i]\,\E_{\rho_1}[\gamma^j]$,
because $\alpha \neq \gamma$. For the last claim, induction on $n$ with $i = 1$
and $j = n - 1$ gives $a_n = a_1^n$ for every $n$.
\end{proof}

Theorem~\ref{thm:substitution} assumes that every $\gamma_i$ differs from
$\beta$. This assumption cannot be dropped. Take the expression 
$t(\gamma, x) = \E[\gamma x]$ and put
$\gamma = \beta$. Then $\den{t}{\rho}(\beta, \alpha) = a_1^2$, while
$\den{t}{\rho}(\beta, \beta) = a_2$. These differ when the variance is nonzero.
Theorem~\ref{thm:uncorrelation} shows where the assumption enters the proof. If
$\gamma_i = \beta$, the relation $R_U$ must relate $\alpha$ to $\beta$ and
$\gamma_i$ to $\beta$. So $R_U$ relates two distinct umbrae, $\alpha$ and
$\gamma_i$, to the same umbra $\beta$. Let $R_K$ be equality on the scalars $k$,
and suppose there is an $R_A$ such that $R_K$, $R_U$, $R_A$ respect the moment
dictionary. Then by Theorem~\ref{thm:uncorrelation}, $a_i a_j = a_{i+j}$ for all
$i, j \geq 1$, that is, $a_n = a_1^n$ for every $n$ which determines the moments
of the umbra completely and uniquely in terms of $a_1$. For an umbra
which is not of this form no such $R_A$ exists.

Similar observations exist in the literature. A product of moments $a_i a_j$ is
in general not $\E[\alpha^i \alpha^j]$. It needs two distinct similar umbrae
\cite[\S2]{DiNardoSenatoKstat}, \cite[\S2.1]{DiNardoNiederhausenSenato}. We already made use 
of it while proving proposition \ref{prop:euler}.

Now let us continue milking theorem \ref{thm:free}.
An expression in two umbrae is a closed term of the type
\begin{equation}
\label{eq:twotype}
  \forall K\,U\,A.\ \Umbral(K, U, A) \to U \to U \to K .
\end{equation}
Its free theorem is Theorem~\ref{thm:free} with $m = 2$: if
$u \mathrel{R_U} v$ and $u' \mathrel{R_U} v'$, then
$\den{t}{d_1}(u, u') \mathrel{R_K} \den{t}{d_2}(v, v')$.

\begin{theorem}[Similar Umbrae Are Exchangeable]
\label{thm:exchange}
Let $\alpha \equiv \beta$ be similar umbrae. Then
$\den{t}{\rho}(\alpha, \beta) = \den{t}{\rho}(\beta, \alpha)$ for every expression
$t$ of type~\eqref{eq:twotype}.
\end{theorem}

\begin{proof}
Take the relations of the proof of Theorem~\ref{thm:substitution}, built from the
transposition $\theta$ exchanging $\alpha$ and $\beta$. The graph of $\theta$
relates $\alpha$ to $\beta$ and $\beta$ to $\alpha$, so the free theorem gives
$\den{t}{\rho}(\alpha, \beta) = \den{t}{\rho}(\beta, \alpha)$.

If a third umbra $\delta \equiv \alpha$, distinct from $\alpha$ and $\beta$, is
available, the theorem also follows from three applications of the substitution
principle:
\[
  \den{t}{\rho}(\alpha, \beta) = \den{t}{\rho}(\delta, \beta)
    = \den{t}{\rho}(\delta, \alpha) = \den{t}{\rho}(\beta, \alpha) .
\]
\end{proof}

For instance, $\E[\alpha^2\beta] = a_2 b_1$ and $\E[\alpha\beta^2] = a_1 b_2$
agree when $\alpha \equiv \beta$, and differ in general otherwise. It is worth
noting that under the probability theoretic interpretation of umbral calculus the
probabilistic notion of exchangeability corresponds this notion exchangeability. Recall 
that a sequence $(X_1, \dots, X_n)$ of random variables is \emph{exchangeable} if
$(X_{\pi(1)}, \dots, X_{\pi(n)})$ has the same joint distribution as
$(X_1, \dots, X_n)$ for every permutation $\pi$ \cite[\S1]{Aldous1985}. For
independent, identically distributed random variables $X$ and $Y$, the pairs
$(X, Y)$ and $(Y, X)$ have the same joint distribution. Rota and Taylor use the word
\emph{exchangeable} for umbrae with the same moments \cite[\S4]{Gessel2003}.

The final theorem of this section is a restriction on the values of expressions.
Let $k_\rho$ be the subring of $k$ generated by all the moments $\rho(\alpha)_i$.

\begin{theorem}[Integrality]
\label{thm:integral}
For every expression $t$ and every umbra $\alpha$, the value
$\den{t}{\rho}(\alpha)$ lies in $k_\rho$.
\end{theorem}

\begin{proof}
Take $R_K$ to relate each element of $k_\rho$ to itself. So $R_K$ is a subset of 
equality on $k$. Take $R_U$ to be equality and $R_A$ to relate each polynomial
with all coefficients in $k_\rho$ to itself. These respect the moment dictionary
of $\rho$ and itself, that is, Definition~\ref{def:respects} holds with $d_1 = d_2$
the moment dictionary of $\rho$. A coefficient of a product is a sum of products of coefficients.
Constants and variables have coefficients $0$, $1$ or an element of $k_\rho$. For
the evaluation, $\E_\rho[p]$ is a sum of coefficients of $p$ times products of
moments, and so lies in $k_\rho$. Since $\alpha \mathrel{R_U} \alpha$, the free
theorem gives $\den{t}{\rho}(\alpha) \mathrel{R_K} \den{t}{\rho}(\alpha)$, that
is, $\den{t}{\rho}(\alpha) \in k_\rho$.
\end{proof}

\section{Naturality versus Parametricity}
\label{sec:naturality}

This is all fine but you probably noticed that we developed a lot of machinery 
and proved modest theorems. For instance, we have not used the full relational
tools and used mostly graphs of functions. Unfortunately it is not a coincidence. 
There is not much one can obtain using proper relations. The last section of this
chapter will make this claim precise.

\begin{definition}
\label{def:hom}
A \emph{homomorphism} of dictionaries $h : d_1 \to d_2$ is a triple of maps
$h_K : \sK_1 \to \sK_2$, $h_U : \sU_1 \to \sU_2$ and $h_A : \sA_1 \to \sA_2$ that commutes
with all methods: $h_K$ and $h_A$ preserve the ring operations,
$h_A \circ \mathrm{const} = \mathrm{const} \circ h_K$,
$h_A \circ \mathrm{inj} = \mathrm{inj} \circ h_U$, and
$h_K \circ \E = \E \circ h_A$.
\end{definition}

\begin{definition}
\label{def:natural}
A \emph{class function on dictionaries} is a class $F$ of pairs $(d, F_d)$, one
for each dictionary $d$, where $F_d : \sU \to \sK$ is a function and $\sK$, $\sU$
are the carriers of $d$.
\begin{itemize}
\item $F$ is \emph{parametric} if, for all dictionaries $d_1$ and $d_2$, all
  relations $R_K$, $R_U$, $R_A$ that respect $d_1$ and $d_2$, and all
  $u \mathrel{R_U} v$, we have $F_{d_1}(u) \mathrel{R_K} F_{d_2}(v)$.
\item $F$ is \emph{natural} if $h_K \circ F_{d_1} = F_{d_2} \circ h_U$ for every
  homomorphism $h : d_1 \to d_2$.
\end{itemize}
\end{definition}

\begin{theorem}
\label{thm:natural}
A class function on dictionaries is parametric if and only if it is natural.
\end{theorem}

\begin{proof}
The graphs of the three maps of a homomorphism respect $d_1$ and $d_2$, which
gives one direction. Conversely, let $R_K$, $R_U$, $R_A$ respect $d_1$ and
$d_2$. Every method
takes arguments of the sorts $K$, $U$, $A$, so the related pairs are closed under
the methods. Thus the pairs in $R_K$, in $R_U$ and in $R_A$ form a dictionary $d$,
with each method acting componentwise. The two projections out of $d$ are
homomorphisms to $d_1$ and $d_2$. Suppose that $u \mathrel{R_U} v$. Then $(u, v)$
is an umbra of $d$, and $F_d(u, v)$ is a pair in $R_K$. By naturality along the
two projections, its components are $F_{d_1}(u)$ and $F_{d_2}(v)$.
\end{proof}

A note on formalization in case you are into that kind of stuff: The theorem
quantifies over class functions. In Lean it is not an issue but in ZFC you actually have a schema.

\begin{corollary}
\label{cor:natural}
For every expression $t$, the class function $d \mapsto \den{t}{d}$ is natural.
\end{corollary}

\begin{proof}
By Theorem~\ref{thm:free} with $m = 1$ it is parametric.
\end{proof}

The results of Section~\ref{sec:free} that are instances of
Theorem~\ref{thm:free} are naturality along homomorphisms. For example:
\begin{itemize}
\item The substitution principle (Theorem~\ref{thm:substitution}) is naturality along the homomorphism
  $(\mathrm{id}, \theta, \theta_*)$ induced by the transposition of two similar
  umbrae; Lemma~\ref{lem:rename} is the condition for the evaluation.
\item Integrality (Theorem~\ref{thm:integral}) is naturality along the inclusion of the moment dictionary over
  $k_\rho$ into the moment dictionary over $k$.
\end{itemize}

The last theorem of this section is about a soundness theorem for the probabilistic
interpretation of umbral calculus.

\begin{definition}
\label{def:probdict}
Let $(\Omega, \mu)$ be a probability space and $X = (X_\alpha)_{\alpha \in
\mathcal{A}}$ a family of real functions on $\Omega$. The \emph{probabilistic
dictionary} of $X$ has $\sK = \RR$, $\sU = \mathcal{A}$ and $\sA$ the set of functions
$\Omega \to \RR$, the pointwise ring operations on $\sA$, the constant functions as
$\mathrm{const}$, $\mathrm{inj}(\alpha) = X_\alpha$, and the integral
$\E(f) = \int f\, d\mu$ as evaluation.
\end{definition}

The following theorem compares two dictionaries: the moment dictionary
of $\rho$ and the probabilistic dictionary of $X$. Both have scalars
$\RR$ and umbrae $\mathcal{A}$. In the moment dictionary of $\rho$, the umbral
polynomials are formal polynomials in $\RR[\mathcal{A}]$. In the probabilistic
dictionary, they are functions on $\Omega$. Substituting the random variables,
$p \mapsto p(X)$, turns a formal polynomial into a function. The theorem says
that an expression has the same value in both dictionaries.

\begin{theorem}[Probabilistic Soundness]
\label{thm:probabilistic}
Let $X$ be a model of an assignment $\rho$ over $\RR$
(Definition~\ref{def:probmodel}). Then, for every expression
$t$ and every name $\alpha$, the value of $t$ at $\alpha$ in the probabilistic
dictionary of $X$ is $\den{t}{\rho}(\alpha)$.
\end{theorem}

\begin{proof}
Let $h_A(p)$ be the function $\omega \mapsto p(X(\omega))$ obtained by
substituting the random variables into the polynomial $p$. Together with the
identities $h_K$ on $\RR$ and $h_U$ on the names it commutes with the ring
operations, with $\mathrm{const}$ and with $\mathrm{inj}$, because substitution
is a ring homomorphism. For the evaluation, $\int h_A(p)\, d\mu = \E_X[p] =
\E_\rho[p]$ by Proposition~\ref{prop:probmodel}. Hence
$h$ is a homomorphism from the moment dictionary of $\rho$ to the probabilistic
dictionary of $X$, and Corollary~\ref{cor:natural} gives the claim.
\end{proof}


\chapter{Umbral Calculus with Resource Allocation}
\label{ch:full}

The previous chapter ended in a somewhat anticlimactic way, by proving that
full relational parametricity was not really necessary for the restricted language. This is 
not really surprising because that language had very little expressive power. For instance
it did not support conjuring a new umbra which we saw to be useful in Proposition~\ref{prop:euler}.
In this chapter we will work with a more expressive language and prove free theorems about it which do 
require relations.


\section{The Language with Allocation}
\label{sec:alloc-language}

The \emph{full language} extends the restricted language of the previous chapter. Its types gain 
products, a type $\Nat$ of natural numbers and a type former $\T$:
\[
  \tau \;::=\; X \mid \tau \to \tau \mid \tau \times \tau \mid \Nat \mid \T\tau
    \mid \forall X.\, \tau .
\]
The terms gain pairs $\langle M, N \rangle$ with projections $\pi_1$ and $\pi_2$,
$0 : \Nat$, $\mathrm{succ} : \Nat \to \Nat$, an iterator
$\mathrm{iter}_\tau : \Nat \to (\tau \to \tau) \to \tau \to \tau$ at every type
$\tau$, and the monadic operations
\[
  \mathrm{return} : \tau \to \T\tau ,
  \qquad
  \mathrm{bind} : \T\sigma \to (\sigma \to \T\tau) \to \T\tau ,
\]
written $x \leftarrow M;\, N$ for $\mathrm{bind}\, M\, (\lambda x.\, N)$. The
dictionary gains one method, after the thirteen of Definition~\ref{def:class},
\[
  \mk : (\Nat \to K) \to \T U ,
\]
the allocation of a fresh umbra whose moments are the given sequence. Again its
zeroth entry is never used. From now on $\Umbral(K, U, A) \to \tau$ abbreviates
the curried function type that takes these fourteen methods. An 
\emph{expression of the full language in $m$ umbrae} is a closed term of type
$\forall K\,U\,A.\ \Umbral(K, U, A) \to U \to \dots \to U \to \T K$.

This new language is really more expressive. We can define Rota and Taylor's auxiliary umbrae, 
such as the dot product and the composition of umbrae using $\mk$. But first we need 
some preparation.

\begin{definition}
\label{def:bell}
For a sequence $a = (a_1, a_2, \dots)$ in a commutative ring and $m, j \in \NN$,
the \emph{partial Bell polynomial} $B_{m,j}(a)$ is the sum, over the partitions
$\pi$ of $\{1, \dots, m\}$ into $j$ parts, of $\prod_{P \in \pi} a_{|P|}$, where
$|P|$ is the number of elements of the part $P$.
The \emph{falling factorial} $(n)_j$ is $n(n-1)\cdots(n-j+1)$.
\end{definition}

The partial Bell polynomials satisfy
$B_{m,j} = \sum_{i=1}^{m-j+1} \binom{m-1}{i-1} a_i\, B_{m-i,\,j-1}$, with
$B_{0,0} = 1$ and $B_{m,0} = B_{0,j} = 0$ otherwise. Indeed, in a partition of
$\{1, \dots, m\}$ into $j$ parts, let $i$ be the size of the part containing
$m$. There are $\binom{m-1}{i-1}$ choices for its other elements, and the
remaining $m - i$ elements form a partition into $j - 1$ parts. This recurrence
is a term, by iteration on $\Nat$ at the type $\Nat \to \Nat \to K$ of tables of
scalars. The power $\alpha^n$ for a variable $n$ is a term. So is the moment
sequence $n \mapsto \E[\alpha^n]$, and so are the factorial moments
$\E[(\beta)_j]$.

In the terms below we write $\alpha$ for $\mathrm{inj}\,\alpha$ wherever a
polynomial is expected, as in $\alpha_1 + \dots + \alpha_n$ and $\E[\alpha^i]$.

\begin{definition}
\label{def:auxiliary}
The \emph{auxiliary umbrae} are the following terms, where $u$, $\alpha$ and
$\beta$ have type $U$ and $n$ has type $\Nat$.
\begin{enumerate}
\item The \emph{fresh copy} is the term
  \[
    \mathrm{new} = \lambda u.\ \mk\,\bigl(\lambda i.\ \E\bigl(\mathrm{iter}_A\, i\,
      (\lambda q.\ q \cdot_A \mathrm{inj}\, u)\, 1_A\bigr)\bigr) ,
  \]
  which allocates an umbra whose $i$-th moment is $\E[u^i]$.
\item The \emph{constructed dot product} is the term $\mathrm{dot}_{\mathrm{c}}$
  defined, by iteration on $n$, by
  \[
    \mathrm{dot}_{\mathrm{c}}\, n\, u = \alpha_1 \leftarrow \mathrm{new}\, u;\ \dots;\
    \alpha_n \leftarrow \mathrm{new}\, u;\
    \mathrm{return}\,(\alpha_1 + \dots + \alpha_n) .
  \]
\item The \emph{symbolic dot product}, after Rota and Taylor
  \cite{RotaTaylor1994} and Di Nardo and Senato \cite[\S2.2]{DiNardoSenato2001},
  is the term $\mathrm{dot}_{\mathrm{s}}$ defined by
  \[
    \mathrm{dot}_{\mathrm{s}}\, n\, u = \gamma \leftarrow \mk\,\Bigl(m \mapsto
    \sum_{j} (n)_j\, B_{m,j}(\E[u], \E[u^2], \dots)\Bigr);\ \mathrm{return}\,\gamma .
  \]
\item The \emph{composition} of $\beta$ and $\alpha$, after Di Nardo and Senato
  \cite[\S2.2]{DiNardoSenato2001}, is the term
  \[
    \beta.\alpha = \mk\,\Bigl(m \mapsto \sum_j \E[(\beta)_j]\, B_{m,j}(\E[\alpha],
    \E[\alpha^2], \dots)\Bigr) .
  \]
\item The \emph{singleton umbra} is the term $\chi = \mk\,(m \mapsto \chi_m)$,
  where $\chi_1 = 1$ and $\chi_m = 0$ for $m \neq 1$.
\item The \emph{cumulant umbra} of $\alpha$, after Di Nardo and Senato
  \cite{DiNardoSenato2006}, is the term $\chi.\alpha = c \leftarrow \chi;\ c.\alpha$,
  where $\chi$ is the term of (5).
\end{enumerate}
\end{definition}

The compositional inverse is a term as well but it is not straightforward to prove. 
Let $\alpha$ have moments $a$ with $a_1 = 1$, and let $f(t) = \sum_{m \geq 0} a_m t^m / m!$. 
Di Nardo and Senato's compositional inverse of $\alpha$ is an umbra $\gamma$ whose moments $g$ have
generating function $g(t) = \sum_{m \geq 0} g_m t^m / m!$ with
$g[f(t) - 1] = f[g(t) - 1] = 1 + t$. Since
$\sum_{m} B_{m,i}(a)\, t^m / m! = (f(t) - 1)^i / i!$, the first equation says
$\sum_{i \leq m} g_i\, B_{m,i}(a) = \chi_m$ for every $m$, and in the same way the
second equation says $\sum_{i \leq m} a_i\, B_{m,i}(g) = \chi_m$ for every
$m \geq 1$.

\begin{definition}
\label{def:inverse}
Let $a$ be a sequence with $a_1 = 1$. The \emph{inversion step} at $a$ is the
map $\Phi$ on sequences given by
\[
  \Phi(s)_j = \chi_j - \sum_{i \leq j} s_i\, B_{j,i}(a) + s_j .
\]
The \emph{inverse moments} are $g_m = \Phi^m(\chi)_m$, and the term
$\mathrm{inverse}\,u = \mk\,(m \mapsto \Phi^m(\chi)_m)$, with $a$ the moments
of $u$, is the \emph{compositional inverse}.
\end{definition}

\begin{proposition}
\label{prop:inverse}
Let $a$ be a sequence with $a_1 = 1$, and let $g$ be its inverse moments. Then
$\sum_{i \leq m} g_i\, B_{m,i}(a) = \chi_m$ for every $m$. Moreover every sequence
$g'$ that satisfies these equations for $m \geq 1$ agrees with $g$ from index~$1$
on.
\end{proposition}

\begin{proof}
The diagonal of the partial Bell polynomials is $B_{j,j}(a) = a_1^j = 1$. Hence
$\Phi(s)_j = \chi_j - \sum_{i < j} s_i\, B_{j,i}(a)$, which depends only on the
entries of $s$ below $j$. By induction on $j$, the $j$-th entry of $\Phi^r(\chi)$
does not depend on $r$ once $r \geq j$, so $g_j = \chi_j - \sum_{i < j} g_i\,
B_{j,i}(a)$. Adding $g_j = g_j\, B_{j,j}(a)$ to both sides gives the equation.
For uniqueness, the equation at $m \geq 1$ determines $g'_m$ from
$g'_1, \dots, g'_{m-1}$ in the same way, and $B_{m,0}(a) = 0$ for $m \geq 1$, so
$g'_0$ does not enter.
\end{proof}

This implies that the compositional inverse is definable in the full language. Indeed, the
partial Bell polynomials are a term by the recurrence above, so the inversion
step $\Phi$ is a term at the type $(\Nat \to K) \to (\Nat \to K)$ of maps on
sequences. Its iterates $\Phi^m(\chi)$ are a term by iteration on $\Nat$ at that
type. Applying $\mk$ to $m \mapsto \Phi^m(\chi)_m$ gives the inverse. 

\section{A Kripke Semantics}
\label{sec:worlds}

Previous section was all syntax. Now we will define a semantics for the language
using Kripke relations.

\begin{definition}
\label{def:world}
A \emph{world} is a pair $w = (n, \rho_w)$ of a natural number $n$ and a
function $\rho_w : [n] \to k^{\NN}$, the moments, where
$[n] = \{0, \dots, n - 1\}$. Its set of \emph{names} is $|w| = [n]$. Its
\emph{evaluation} $\E_w$ is the evaluation of the moment model of $\rho_w$ on
$k[|w|]$. A \emph{morphism} $\kappa : w \to v$ is an injective map
$\kappa : |w| \to |v|$ with $\rho_v(\kappa x)_j = \rho_w(x)_j$ for all $x$ and
all $j \geq 1$. Worlds and morphisms form a category $\W$. The \emph{extension}
of $w = (n, \rho_w)$ by a world $\Delta = (m, \rho_\Delta)$, called a
\emph{block}, is the world $w \oplus \Delta = (n + m, \rho)$ with
$\rho(x) = \rho_w(x)$ for $x < n$ and $\rho(n + y) = \rho_\Delta(y)$ for $y < m$.
\end{definition}

The morphisms of $\W$ are the moment-preserving renamings of
Lemma~\ref{lem:rename}, allowed to change the set of names. Extension by a fixed 
world is functorial. This means a morphism $\kappa : w \to v$ induces
$\kappa \oplus \Delta : w \oplus \Delta \to v \oplus \Delta$.

\begin{definition}
\label{def:kripke-semantics}
A \emph{$\W$-set} is a functor $\mathcal{F} : \W \to \mathbf{Set}$. The types
denote $\W$-sets as follows. The type $\Nat$ denotes the constant $\W$-set
$\NN$, and products are pointwise. The \emph{function
space} $\mathcal{F} \Rightarrow \mathcal{F}'$ is the $\W$-set
\[
  (\mathcal{F} \Rightarrow \mathcal{F}')(w)
    = \prod_{\kappa : w \to v} \bigl(\mathcal{F}(v) \to \mathcal{F}'(v)\bigr),
  \qquad (\kappa' \cdot f)_{\kappa} = f_{\kappa \circ \kappa'} ,
\]
for morphisms $\kappa' : w \to w'$ and $\kappa : w' \to v$. For a $\W$-set
$\mathcal{F}$, let $\T\mathcal{F}$ be the $\W$-set with
\[
  (\T\mathcal{F})(w) = \{(\Delta, s) : \Delta \text{ a world},\
    s \in \mathcal{F}(w \oplus \Delta)\},
  \qquad \kappa \cdot (\Delta, s) = (\Delta, (\kappa \oplus \Delta) \cdot s) .
\]
The \emph{allocation monad} is $\T$, a monad on the category of $\W$-sets. Its
unit is $s \mapsto (\varnothing, s)$, where $\varnothing = (0, \emptyset)$ is the
world with no names, and its multiplication takes $(\Delta, (\Delta', s))$ to
$(\Delta \oplus \Delta', s)$. For a world $w$ and a block $\Delta$, write
$\mathrm{inl} : w \to w \oplus \Delta$ for the inclusion of names.
\end{definition}

The semantics follows Definition~\ref{def:semantics} with $\W$-sets in place of
sets. It is the functor-category semantics of local names of Pitts and Stark
\cite{PittsStark1993} and of O'Hearn and Tennent \cite{OHearnTennent1995}, with
the allocation monad above. I am reproducing it here to make the notes self contained.

\begin{definition}
\label{def:kripke-terms}
An \emph{interpretation} is a function $I$ from the base sorts to $\W$-sets.
The simple types denote $\W$-sets by $\den{b}{} = I(b)$ and
$\den{\sigma \to \sigma'}{} = \den{\sigma}{} \Rightarrow \den{\sigma'}{}$. For a
type environment $\eta$, as in Definition~\ref{def:semantics}, the types denote
$\W$-sets by
\begin{gather*}
  \den{X}{\eta} = \den{\eta(X)}{}, \qquad
  \den{\tau \to \tau'}{\eta} = \den{\tau}{\eta} \Rightarrow \den{\tau'}{\eta}, \qquad
  \den{\tau \times \tau'}{\eta}(w) = \den{\tau}{\eta}(w) \times \den{\tau'}{\eta}(w), \\
  \den{\Nat}{\eta}(w) = \NN, \qquad
  \den{\T\tau}{\eta} = \T\den{\tau}{\eta}, \qquad
  \den{\forall X.\, \tau}{\eta}(w) = \prod_{\sigma \text{ simple}}
    \den{\tau}{\eta[X \mapsto \sigma]}(w),
\end{gather*}
and a context $\Gamma$ denotes the pointwise product $\den{\Gamma}{\eta}$ of the
denotations of its types. A term $\Gamma \vdash M : \tau$ denotes, at each world
$w$, a function $\den{M}{w} : \den{\Gamma}{\eta}(w) \to \den{\tau}{\eta}(w)$,
defined by induction on $M$:
\begin{itemize}
\item $\den{x}{w}(\gamma) = \gamma(x)$;
\item $\den{\lambda x.\, M}{w}(\gamma)$ is the family whose component at
  $\kappa : w \to v$ is $a \mapsto \den{M}{v}(\kappa \cdot \gamma, a)$;
\item $\den{M\,N}{w}(\gamma) = \den{M}{w}(\gamma)_{\mathrm{id}_w}\bigl(\den{N}{w}(\gamma)\bigr)$;
\item pairs and projections act componentwise, $\den{0}{w} = 0$ and
  $\den{\mathrm{succ}\,M}{w}(\gamma) = \den{M}{w}(\gamma) + 1$;
\item $\den{\mathrm{iter}\,M\,F\,N}{w}(\gamma) = f^n(x)$, where
  $n = \den{M}{w}(\gamma)$, $f = \den{F}{w}(\gamma)_{\mathrm{id}_w}$ and
  $x = \den{N}{w}(\gamma)$;
\item $\den{\mathrm{return}\,M}{w}(\gamma) = (\varnothing, \den{M}{w}(\gamma))$;
\item $\den{\mathrm{bind}\,M\,F}{w}(\gamma) = (\Delta \oplus \Delta', s')$, where
  $(\Delta, s) = \den{M}{w}(\gamma)$ and
  $(\Delta', s') = \den{F}{w}(\gamma)_{\mathrm{inl}}(s)$ for the inclusion
  $\mathrm{inl} : w \to w \oplus \Delta$;
\item $\den{\Lambda X.\, M}{w}$ is the family
  $\bigl(\den{M}{w}\bigr)_\sigma$, the denotation of $M$ under
  $\eta[X \mapsto \sigma]$ for each simple type $\sigma$.
\end{itemize}
\end{definition}

In the clause for $\mathrm{bind}$, the continuation runs in the extended world
$w \oplus \Delta$, and $(w \oplus \Delta) \oplus \Delta' = w \oplus (\Delta \oplus
\Delta')$. The formal version is \texttt{Core/Kripke.lean} (\texttt{QTyD},
\texttt{eval}); see Chapter~\ref{ch:formal}.

\begin{definition}
\label{def:momentdict}
The \emph{moment interpretation} is the interpretation with three $\W$-sets,
\[
  \sK(w) = k , \qquad \sU(w) = |w| , \qquad \sA(w) = k[|w|] .
\]
A morphism $\kappa$ acts on $\sK$ as the identity, on $\sU$ as $\kappa$, and on
$\sA$ by renaming, written $\kappa_*$.

The \emph{moment dictionary} at $w$ is a tuple with one element for each method.
The element for a method of type $\mu$ lies in $\den{\mu}{}(w)$. An element of a
function type is a family indexed by morphisms $\kappa : w \to v$. The elements
are as follows.
\begin{itemize}
\item The constants are $0$ and $1$ of $k$ and of $k[|w|]$.
\item Let $\circ$ be a binary ring operation. It takes $x$ at $\kappa : w \to v$
  and $y$ at $\kappa' : v \to v'$ to $\kappa'_*(x) \circ y$. On $k$, $\kappa'_*$ is
  the identity.
\item Negation at $\kappa : w \to v$ is negation at $v$.
\item At $\kappa : w \to v$, $\mathrm{const}$, $\mathrm{inj}$ and $\E$ are those of
  the moment model of $\rho_v$.
\item At $\kappa : w \to v$, $\mk$ takes $f \in \den{\Nat \to K}{}(v)$ to
  $(\Delta_f, \bullet) \in (\T \sU)(v)$. Here $\Delta_f$ is the block with one
  name and moments $f_{\mathrm{id}_v}$, and $\bullet$ is that name in
  $v \oplus \Delta_f$. If $v$ has $n$ names, then $\bullet = n$.
\end{itemize}
\end{definition}

The allocated name lies in the block, so it is distinct from every name of $w$,
and the evaluation of $w \oplus \Delta_f$ treats it as uncorrelated with all of
them. No axiom is added for freshness: uncorrelation of distinct names is the
moment model itself.

\begin{definition}
\label{def:worldrel}
Let $w_1$ and $w_2$ be worlds. A \emph{world relation}
$\varphi : w_1 \nrightarrow w_2$ is a $k$-algebra homomorphism
$\varphi : k[|w_1|] \to k[|w_2|]$ with $\E_{w_2} \circ \varphi = \E_{w_1}$. Its
\emph{name part} is the relation
$N_\varphi = \{(x, y) : \varphi(x) = y\} \subseteq |w_1| \times |w_2|$.
\end{definition}

Every morphism $\kappa$ of $\W$ induces a world relation $\kappa_*$, the renaming
of variables, by Lemma~\ref{lem:rename}. Its name part is the graph of $\kappa$.
The converse fails: the algebra
map $\gamma \mapsto \alpha_1 + \dots + \alpha_n$ of section~\ref{sec:repind} is a
world relation whose name part does not contain $\gamma$ at all.

\begin{lemma}[Extension]
\label{lem:extension}
Let $\varphi : w_1 \nrightarrow w_2$ and $\psi : \Delta_1 \nrightarrow \Delta_2$
be world relations. Then
$\varphi \otimes \psi : k[|w_1 \oplus \Delta_1|] \to k[|w_2 \oplus \Delta_2|]$,
the algebra map that is $\varphi$ on the names of $w_1$ and $\psi$ on those of
$\Delta_1$, is a world relation
$w_1 \oplus \Delta_1 \nrightarrow w_2 \oplus \Delta_2$.
\end{lemma}

\begin{proof}
By linearity it suffices to treat a product $p\,q$ with $p \in k[|w_1|]$ and
$q \in k[|\Delta_1|]$, since such products span $k[|w_1 \oplus \Delta_1|]$. The
image is $\varphi(p)\,\psi(q)$, a product of a polynomial in the names of $w_2$
and a polynomial in the names of $\Delta_2$. These two sets of names are disjoint,
so uncorrelation factors both evaluations:
\[
  \E\bigl[\varphi(p)\,\psi(q)\bigr] = \E[\varphi(p)]\,\E[\psi(q)]
    = \E[p]\,\E[q] = \E[p\,q] . \qedhere
\]
\end{proof}

\begin{definition}
\label{def:future}
A \emph{future} of a world relation $\varphi : w_1 \nrightarrow w_2$ is a triple
$(\kappa_1, \kappa_2, \varphi')$ of morphisms $\kappa_i : w_i \to v_i$ and a world
relation $\varphi' : v_1 \nrightarrow v_2$ with
$\varphi' \circ (\kappa_1)_* = (\kappa_2)_* \circ \varphi$.
\end{definition}

Futures compose, and $(\mathrm{id}, \mathrm{id}, \varphi)$ is a future of
$\varphi$. For blocks $\Delta_1$, $\Delta_2$ and a world relation
$\psi : \Delta_1 \nrightarrow \Delta_2$, the inclusions together with
$\varphi \otimes \psi$ form a future of $\varphi$, by Lemma~\ref{lem:extension}.

\begin{definition}
\label{def:kripke-logrel}
Let $\varphi : w_1 \nrightarrow w_2$ be a world relation. The \emph{logical
relation} at a type $\tau$ is the relation
$\rel{\tau}{\varphi} \subseteq \den{\tau}{}(w_1) \times \den{\tau}{}(w_2)$
defined by induction on $\tau$ as follows. Type variables carry relations as in
Definition~\ref{def:logrel}. For the moment dictionary they are instantiated at
the three sorts below.
\begin{itemize}
\item $\rel{K}{\varphi}$ and $\rel{\Nat}{\varphi}$ are equality,
  $\rel{U}{\varphi} = N_\varphi$, and $\rel{A}{\varphi}$ is the graph of
  $\varphi$.
\item $(x, y) \mathrel{\rel{\tau \times \tau'}{\varphi}} (x', y')$ if and only if
  $x \mathrel{\rel{\tau}{\varphi}} x'$ and $y \mathrel{\rel{\tau'}{\varphi}} y'$.
\item $f \mathrel{\rel{\tau \to \tau'}{\varphi}} g$ if and only if
  $f_{\kappa_1}(x) \mathrel{\rel{\tau'}{\varphi'}} g_{\kappa_2}(y)$ for every
  future $(\kappa_1, \kappa_2, \varphi')$ of $\varphi$ and all
  $x \mathrel{\rel{\tau}{\varphi'}} y$.
\item $(\Delta_1, s_1) \mathrel{\rel{\T\tau}{\varphi}} (\Delta_2, s_2)$ if and
  only if there is a world relation $\psi : \Delta_1 \nrightarrow \Delta_2$ with
  $s_1 \mathrel{\rel{\tau}{\varphi \otimes \psi}} s_2$.
\end{itemize}
\end{definition}

The clause for $\T$ is very much like the Kripke lifting of Pitts and Stark
\cite{PittsStark1993}. In the version here the relation between the two blocks is an
arbitrary world relation as opposed to a bijection of names.

\begin{lemma}[Monotonicity]
\label{lem:monotone}
If $x \mathrel{\rel{\tau}{\varphi}} y$ and $(\kappa_1, \kappa_2, \varphi')$ is a
future of $\varphi$, then
$\kappa_1 \cdot x \mathrel{\rel{\tau}{\varphi'}} \kappa_2 \cdot y$.
\end{lemma}

\begin{proof}
By induction on $\tau$. At $K$ and $\Nat$ there is nothing to prove. At $A$,
$\varphi'(\kappa_1 x) = \kappa_2(\varphi x) = \kappa_2 y$ by the defining equation
of a future, and the case of $U$ is the same equation at a variable. At a function
type it follows because futures compose. At $\T\tau$, let $\psi$ witness the
relation of $(\Delta_1, s_1)$ and $(\Delta_2, s_2)$. Then
$(\kappa_1 \oplus \Delta_1, \kappa_2 \oplus \Delta_2, \varphi' \otimes \psi)$ is a
future of $\varphi \otimes \psi$: both sides of the defining equation agree on the
names of $w_1$ by hypothesis and on those of $\Delta_1$ by construction, and
$\varphi' \otimes \psi$ is a world relation by Lemma~\ref{lem:extension}. The
induction hypothesis at $\tau$ concludes.
\end{proof}

\begin{lemma}
\label{lem:dictrel}
At every world relation $\varphi$, the moment dictionary is related to itself:
each method is related to itself by the logical relation at its type.
\end{lemma}

\begin{proof}
The ring operations of $A$ are related because $\varphi$ is a ring homomorphism,
and $\mathrm{const}$ because it is a $k$-algebra homomorphism. The method
$\mathrm{inj}$ sends $x \mathrel{N_\varphi} y$ to $\varphi(x) = y$, which is the
graph of $\varphi$. The evaluation is related because $\varphi$ preserves it.
These statements persist along futures by Lemma~\ref{lem:monotone}. For $\mk$,
let $f \mathrel{\rel{\Nat \to K}{\varphi'}} g$; then $f$ and $g$ agree as
sequences. The two allocations are $(\Delta_f, \bullet)$ and
$(\Delta_g, \bullet)$. Let $\psi : \Delta_f \nrightarrow \Delta_g$ send $\bullet$
to $\bullet$. It preserves the evaluation because the two blocks have the same
moments. The two values are related at $U$, because
$\bullet \mathrel{N_{\varphi' \otimes \psi}} \bullet$.
\end{proof}

And finally\ldots

\begin{theorem}[Abstraction Theorem with Allocation]
\label{thm:abstraction2}
Let $\varphi$ be a world relation, and let the dictionaries and the environments
be related at $\varphi$. Then the denotations of every term are related at its
type.
\end{theorem}

\begin{proof}
By induction on the term, as for Theorem~\ref{thm:abstraction-precise}. We give
the new cases; the formal proof covers all of them (Chapter~\ref{ch:formal}). For
$\lambda$, use the futures in the clause for functions. For a variable that has
been transported to a future, use Lemma~\ref{lem:monotone}. For
$\mathrm{iter}$, use induction on the natural number, which is the same on both
sides. For $\mathrm{return}$, take $\psi$ to be the identity of the empty block.
For $\mathrm{bind}$, let $\psi$ witness the relation of the first computation.
The continuation runs at the future
$(\mathrm{incl}, \mathrm{incl}, \varphi \otimes \psi)$, and its result is related
at some $(\varphi \otimes \psi) \otimes \psi' = \varphi \otimes (\psi \otimes \psi')$.
The witness is therefore $\psi \otimes \psi'$, a world relation by
Lemma~\ref{lem:extension}.
\end{proof}

With Lemma~\ref{lem:dictrel}, the theorem gives the free theorem of the type of
expressions of the full language in $m$ umbrae,
\[
  \forall K\,U\,A.\ \Umbral(K, U, A) \to U \to \dots \to U \to \T K .
\]
For every expression of this type, a world
relation that sends the arguments on one side to those on the other relates the
scalar values, whatever the expression allocates. Its first
instance, at the renaming that exchanges two similar names, is the substitution
principle with allocation.

\begin{proposition}[Conservativity]
\label{prop:conservative}
Let $t$ be an expression of Chapter~\ref{ch:restricted}, read as an expression of the full
language that does not use $\mk$. For every world $w$ and name $u \in |w|$, its
denotation in the full language at the moment dictionary and $w$, applied to $u$,
is its denotation in the restricted language at the moment dictionary of $\rho_w$, applied to $u$.
\end{proposition}

\begin{proof}
A term of the restricted language never allocates, and the only way its denotation consults a
world other than $w$ is through the components $f_\kappa$ of a function. Compare
every function with its component at the identity of $w$. The comparison is a
simulation relation that commutes with every term former of the restricted language, by
induction on the term, and it is equality at the sorts $K$, $U$ and $A$. Two laws
keep it stable: transport along an identity is the identity, and the identity
morphism is a left unit.
\end{proof}

The results of Chapter~\ref{ch:restricted} are statements about these denotations, so
every one of them holds verbatim in the full language. Definition~\ref{def:logrel} relates arbitrary
relations, while Definition~\ref{def:kripke-logrel} relates world relations.

We now compute what the auxiliary umbrae of Definition~\ref{def:auxiliary}
denote at the moment dictionary.

\begin{lemma}[The Moments of $n.\alpha$]
\label{lem:dotmoments}
Let $\alpha_1, \dots, \alpha_n$ be distinct names with the same moments
$a = (a_1, a_2, \dots)$. Then, for every $m$,
\[
  \E\bigl[(\alpha_1 + \dots + \alpha_n)^m\bigr]
    = \sum_{j=0}^{m} (n)_j\, B_{m,j}(a) .
\]
\end{lemma}

\begin{proof}
Expanding the power gives the sum of
$\prod_{r=1}^{m} \alpha_{f(r)}$ over all maps
$f : \{1, \dots, m\} \to \{1, \dots, n\}$. The fibres of $f$ form a partition
$\pi_f$ of $\{1, \dots, m\}$. By uncorrelation the term of $f$ evaluates to
$\prod_{P \in \pi_f} a_{|P|}$, which depends only on $\pi_f$. A partition
with $j$ parts is $\pi_f$ for exactly $(n)_j$ maps $f$, namely the injective
assignments of its parts to indices. Summing over partitions with $j$ parts
gives $(n)_j\, B_{m,j}(a)$.
\end{proof}

\begin{proposition}
\label{prop:auxiliary}
At the moment dictionary of Definition~\ref{def:momentdict}, in a world in which
$u$ has moments $a$, the terms of Definition~\ref{def:auxiliary} denote the
following computations.
\begin{enumerate}
\item $\mathrm{new}\,u$ allocates one name with the moments $a$.
\item $\mathrm{dot}_{\mathrm{c}}\,n\,u$ allocates $n$ names and returns their
  sum, whose moments are the $n$-fold binomial convolution of $a$.
\item $\mathrm{dot}_{\mathrm{s}}\,n\,u$ allocates one name with that sequence of
  moments.
\item $\beta.\alpha$ allocates one name with the moments in
  Definition~\ref{def:auxiliary}(4). These are $\E[q_m(\beta)]$ for the
  polynomial $q_m(x) = \sum_j (x)_j B_{m,j}(a)$, whose value at each $n \in \NN$ is
  the $m$-th moment of $n.\alpha$, which is Rota and Taylor's definition of
  $\beta.\alpha$. For the degenerate umbra with moments $n^i$ they are the moments
  of $n.\alpha$.
\item $\chi.\alpha$ allocates two names, and the $m$-th moment of the one it
  returns is the coefficient of $x$ in $q_m(x)$, the $m$-th cumulant of
  $\alpha$.
\end{enumerate}
\end{proposition}

\begin{proof}
Items (1), (3) and (4) are the definition of $\mk$ at the moment dictionary,
together with Lemma~\ref{lem:dotmoments} for the identification of the sequence
in (3) with the $n$-fold convolution. Item (2) follows by induction on $n$ from
the binomial theorem and (U2), since each new copy is uncorrelated with the
polynomial built so far. In (4) the value of $q_m$ at $n$
is the $m$-th moment of $n.\alpha$ by Lemma~\ref{lem:dotmoments}, and at the
degenerate umbra the evaluation of $q_m(\beta)$ is $q_m(n)$. For (5), the
evaluation of $q_m(\chi)$ keeps only the constant and the linear coefficient of
$q_m$, and the constant coefficient $q_m(0)$ vanishes for $m \geq 1$.
\end{proof}

Note that original similarity was defined between umbrae. We need to 
generalize it because we have now monadic computations producing umbrae.
To be more precise, we need a notion of similarity for values of type $\T U$.

\begin{definition}
\label{def:similar-comp}
Let $c$ and $c'$ be computations of type $\T U$ at the moment dictionary, at
worlds $w$ and $w'$. Each returns a name, of $w \oplus \Delta$ and of
$w' \oplus \Delta'$. Write $c \equiv c'$ if these two names have the same
moments, and $\delta \equiv c$, for an umbra $\delta$, if $\delta$ has the moments
of the name that $c$ returns.
\end{definition}

\begin{proposition}[The auxiliary umbrae respect similarity]
\label{prop:similar}
Let $u \equiv u'$, $\alpha \equiv \alpha'$ and $\beta \equiv \beta'$, where the
umbrae may lie in different worlds.
\begin{enumerate}
\item The results of $\mathrm{new}$, $\mathrm{dot}_{\mathrm{c}}\,n$,
  $\mathrm{dot}_{\mathrm{s}}\,n$, the cumulant umbra and the compositional inverse
  at $u$ and at $u'$ have the same moments.
\item $\beta.\alpha \equiv \beta'.\alpha'$.
\end{enumerate}
No relation between $\beta$ and $\alpha$ is assumed. In particular
$\alpha.\alpha \equiv \alpha'.\alpha'$.
\end{proposition}

\begin{proof}
For (1), Proposition~\ref{prop:auxiliary} and Definition~\ref{def:inverse} give
the moments of each result as a function of the moments of the argument alone.
For $\mathrm{dot}_{\mathrm{c}}$ the result is a polynomial, and its moments are
the values $\E[p^m]$. For (2), expanding the falling factorial gives
$\E[(\beta)_j] = \sum_i s(j, i)\, \E[\beta^i]$, where the $s(j, i)$ are the
Stirling numbers of the first kind. So the factorial moments of $\beta$ are
determined by its moments, and $B_{m,j}(a)$ depends on the moments $a$ of
$\alpha$ alone. Hence $\beta.\alpha$ and $\beta'.\alpha'$ have the same moments.
\end{proof}

\section{Representation Independence}
\label{sec:repind}

This section is about the auxiliary umbrae in this new Kripke-style semantics of 
the language.

\begin{definition}
\label{def:interface}
An \emph{interface} is a type $\tau$ whose free type variables are among $K$,
$U$ and $A$. An \emph{implementation} of $\tau$ is a term of type $\tau$ whose
free variables are among the fourteen methods of $\Umbral(K, U, A)$. A
\emph{context} for $\tau$ is a closed term of type
\[
  \forall K\,U\,A.\ \Umbral(K, U, A) \to \tau \to U \to \T K .
\]
A context $C$ \emph{distinguishes} implementations $i$ and $i'$ if there are a
world $w$ and a name $u \in |w|$ such that, at the moment dictionary and $w$, the
scalar components of $C\,i\,u$ and $C\,i'\,u$ differ.
\end{definition}

\begin{theorem}[Representation Independence of the Dot Product]
\label{thm:repind}
Let $C$ be a context for the interface $\Nat \to U \to \T A$, $w$ a world and $u \in |w|$. Write
$(\Delta_{\mathrm{s}}, c_{\mathrm{s}}) = C\,\mathrm{dot}_{\mathrm{s}}\, u$ and
$(\Delta_{\mathrm{c}}, c_{\mathrm{c}}) = C\,\mathrm{dot}_{\mathrm{c}}\, u$ for
the denotations at the moment dictionary and $w$, with
$c_{\mathrm{s}}, c_{\mathrm{c}} \in k$. Then $c_{\mathrm{s}} = c_{\mathrm{c}}$.
\end{theorem}

\begin{proof}
Apply Theorem~\ref{thm:abstraction2} to $C$ at the identity world relation of $w$.
The dictionary is related to itself by Lemma~\ref{lem:dictrel}, and
$u \mathrel{N_{\mathrm{id}}} u$. It remains to show that
$\mathrm{dot}_{\mathrm{s}} \mathrel{\rel{\Nat \to U \to \T A}{\mathrm{id}}}
\mathrm{dot}_{\mathrm{c}}$.

Let $\varphi : v_1 \nrightarrow v_2$ be a world relation reached by a future, let
$n \in \NN$, and let $u_1 \mathrel{N_\varphi} u_2$. Then $\varphi(u_1) = u_2$, so
$\E[u_1^m] = \E[\varphi(u_1^m)] = \E[u_2^m]$ for every $m$; call this common
sequence $a$. The computation $\mathrm{dot}_{\mathrm{s}}\, n\, u_1$ allocates a
block $\Delta_{\mathrm{s}}$ with one name $\gamma$ and moments
$m \mapsto \sum_j (n)_j B_{m,j}(a)$, and returns $\gamma$. The computation
$\mathrm{dot}_{\mathrm{c}}\, n\, u_2$ allocates a block $\Delta_{\mathrm{c}}$
with $n$ names $\alpha_1, \dots, \alpha_n$, each with moments $a$, and returns
$\alpha_1 + \dots + \alpha_n$. Let $\psi : k[\gamma] \to k[|\Delta_{\mathrm{c}}|]$
be the algebra map with $\psi(\gamma) = \alpha_1 + \dots + \alpha_n$. It
preserves the evaluation, because by Lemma~\ref{lem:dotmoments}
\[
  \E\bigl[\psi(\gamma^m)\bigr] = \E\bigl[(\alpha_1 + \dots + \alpha_n)^m\bigr]
    = \sum_j (n)_j B_{m,j}(a) = \E[\gamma^m] ,
\]
and both sides are linear in the polynomial. So $\psi$ is a world relation, and
$(\varphi \otimes \psi)(\gamma) = \alpha_1 + \dots + \alpha_n$. The two returned
values are therefore related at $A$, and the two computations are related at
$\T A$ with witness $\psi$.

The abstraction theorem gives
$(\Delta_{\mathrm{s}}, c_{\mathrm{s}}) \mathrel{\rel{\T K}{\mathrm{id}}}
(\Delta_{\mathrm{c}}, c_{\mathrm{c}})$. The relation at $K$ is equality, so
$c_{\mathrm{s}} = c_{\mathrm{c}}$.
\end{proof}

The proof has a general form. Two implementations of any interface that are
related at the identity world relation are indistinguishable by every context. The
witness in the proof is the world relation out of a block with one name $\gamma$
that sends $\gamma$ to a polynomial with the moments of $\gamma$. The same
argument gives a second instance.

\begin{corollary}[Only Moments Are Observable]
\label{cor:repind-new}
No context for the interface $U \to \T U$ distinguishes $\mathrm{new}$ from the
implementation $u \mapsto (\_ \leftarrow \mk\,(n \mapsto 0);\ \mathrm{new}\,u)$,
which first allocates an unused umbra.
\end{corollary}

\begin{proof}
As for Theorem~\ref{thm:repind}. The two computations return names with the same
moments. The algebra map that sends the name returned on the left to the name
returned on the right preserves the evaluation, because the two names have the
same moments. So it is a world relation. Moreover the unused name on the right is not in
its image.
\end{proof}

Here comes more computations in the new semantics. At this point the proofs
should feel mechanical.

\begin{lemma}
\label{lem:egf}
Let $a = (a_1, a_2, \dots)$ and $g(t) = \sum_{i \geq 1} a_i t^i / i!$. Then
$\sum_{m} B_{m,j}(a)\, t^m / m! = g(t)^j / j!$. Consequently
\[
  \E\bigl[e^{(p.\alpha) t}\bigr] = \sum_{j \geq 0} \E[(p)_j]\, \frac{g(t)^j}{j!} ,
  \qquad g(t) = \E\bigl[e^{\alpha t}\bigr] - 1 .
\]
\end{lemma}

\begin{proof}
The coefficient of $t^m / m!$ in $g(t)^j$ is the sum, over ordered $j$-tuples
$(i_1, \dots, i_j)$ of positive integers with sum $m$, of
$\binom{m}{i_1, \dots, i_j} \prod_r a_{i_r}$. This counts ordered partitions of
$\{1, \dots, m\}$ into $j$ nonempty parts, weighted by $\prod_P a_{|P|}$. Each
unordered partition into $j$ parts arises from $j!$ orderings. The second
formula follows by summing the moments of $p.\alpha$ against $t^m/m!$; for each
$m$ only $j \leq m$ contribute, because $g$ has no constant term.
\end{proof}

\begin{proposition}[Distributivity]
\label{prop:distributive}
Let $p$ and $q$ be polynomials with no umbra in common, and $\alpha$ an umbra.
Let $\delta \equiv p.\alpha$ and $\delta' \equiv q.\alpha$ be distinct fresh
umbrae. Then $(p + q).\alpha$ has the moments of $\delta + \delta'$.
\end{proposition}

\begin{proof}
In $\ZZ[x, y]$ the falling factorials satisfy the Vandermonde identity
$(x + y)_j = \sum_i \binom{j}{i} (x)_i (y)_{j-i}$. Substitute $p$ and $q$ and
evaluate. Since $p$ and $q$ are uncorrelated,
$\E[(p + q)_j] = \sum_i \binom{j}{i} \E[(p)_i]\, \E[(q)_{j-i}]$. With $y = g(t)$,
Lemma~\ref{lem:egf} gives
\[
  \E\bigl[e^{((p+q).\alpha) t}\bigr]
    = \sum_{j} \E[(p+q)_j] \frac{y^j}{j!}
    = \Bigl(\sum_i \E[(p)_i] \frac{y^i}{i!}\Bigr)\Bigl(\sum_l \E[(q)_l] \frac{y^l}{l!}\Bigr)
    = \E\bigl[e^{\delta t}\bigr]\, \E\bigl[e^{\delta' t}\bigr] .
\]
The last product is $\E[e^{(\delta + \delta') t}]$ by uncorrelation. Comparing
coefficients gives the claim.
\end{proof}

Proposition~\ref{prop:distributive} assumes that $p$ and $q$ have no umbra in
common. We now turn it into a statement about contexts. A context receives its
arguments from outside, and no type can say that two arguments are distinct
umbrae. So the two implementations below first allocate fresh copies $\beta$ and
$\beta'$ of their arguments. Fresh copies are distinct, even when $u = u'$. For
umbrae $u$, $u'$ and $v$, let
\[
  \begin{aligned}
    \mathrm{cs}_{\mathrm{f}}\, u\, u'\, v &= \beta \leftarrow \mathrm{new}\, u;\
      \beta' \leftarrow \mathrm{new}\, u';\ \gamma \leftarrow (\beta + \beta').v;\
      \mathrm{return}\, \gamma , \\
    \mathrm{cs}_{\mathrm{s}}\, u\, u'\, v &= \beta \leftarrow \mathrm{new}\, u;\
      \beta' \leftarrow \mathrm{new}\, u';\ \delta \leftarrow \beta.v;\
      \delta' \leftarrow \beta'.v;\ \mathrm{return}\, (\delta + \delta') ,
  \end{aligned}
\]
be the fused and the split composition of a sum.

\begin{corollary}
\label{cor:repind-compose}
No context for the interface $U \to U \to U \to \T A$ distinguishes
$\mathrm{cs}_{\mathrm{f}}$ from $\mathrm{cs}_{\mathrm{s}}$.
\end{corollary}

\begin{proof}
As for Theorem~\ref{thm:repind}. Related arguments have equal moments. The world
relation on the two blocks fixes $\beta$ and $\beta'$ and sends $\gamma$ to
$\delta + \delta'$. It preserves the evaluation of every monomial
$\beta^i \beta'^j \gamma^l$: all names are distinct on both sides, so both
evaluations factor, and $\E[\gamma^l] = \E[(\delta + \delta')^l]$ by
Proposition~\ref{prop:distributive}.
\end{proof}

\begin{corollary}[Sheffer Identity and Binomial Type at Umbral Arguments]
\label{cor:sheffer}
Let $\alpha$, $u$, $u'$, $v$ be umbrae, let $\beta \equiv u$ and
$\beta' \equiv u'$ be fresh, and let $\delta \equiv (\beta + \beta').v$,
$\delta_1 \equiv \beta.v$ and $\delta_2 \equiv \beta'.v$ be fresh. Then
\[
  \E\bigl[(\alpha + \delta)^n\bigr]
    = \sum_{i=0}^{n} \binom{n}{i} \E\bigl[(\alpha + \delta_1)^i\bigr]\,
      \E\bigl[\delta_2^{\,n-i}\bigr] ,
  \qquad
  \E\bigl[\delta^n\bigr]
    = \sum_{i=0}^{n} \binom{n}{i} \E\bigl[\delta_1^{\,i}\bigr]\, \E\bigl[\delta_2^{\,n-i}\bigr] .
\]
\end{corollary}

\begin{proof}
The umbra $\delta$ is fresh, so it is uncorrelated with $\alpha$, and the
binomial theorem gives
$\E[(\alpha + \delta)^n] = \sum_i \binom{n}{i} \E[\alpha^i]\, \E[\delta^{n-i}]$.
By Proposition~\ref{prop:distributive},
$\E[\delta^l] = \E[(\delta_1 + \delta_2)^l]$, so the right side is
$\E[(\alpha + \delta_1 + \delta_2)^n]$. Expand
$((\alpha + \delta_1) + \delta_2)^n$ by the binomial theorem and evaluate; the
umbra $\delta_2$ is uncorrelated with $\alpha + \delta_1$. The second identity is
the case $\alpha = 0$.
\end{proof}

Write $s_n(\beta) = \E[(\alpha + \beta.v)^n]$ and $p_n(\beta) = \E[(\beta.v)^n]$.
The corollary reads $s_n(\beta + \beta') = \sum_i \binom{n}{i} s_i(\beta)\,
p_{n-i}(\beta')$ and $p_n(\beta + \beta') = \sum_i \binom{n}{i} p_i(\beta)\,
p_{n-i}(\beta')$. These are the Sheffer identity and the binomial type, at umbral
arguments rather than at natural numbers.

Corollary~\ref{cor:inverse-right} below needs associativity of composition. We
state it for moment sequences, which is the form the corollary uses. For
sequences $a$ and $x$, the \emph{composition transform} $C_a(x)$ is the sequence
with $C_a(x)_m = \sum_{i \leq m} x_i\, B_{m,i}(a)$. By
Definition~\ref{def:auxiliary}(4), if $\beta$ has factorial moments $x$ and
$\alpha$ has moments $a$, then $C_a(x)$ is the moment sequence of $\beta.\alpha$.

\begin{proposition}[Associativity of Composition]
\label{prop:assoc}
For all sequences $a$, $b$ and $x$, $C_a(C_b(x)) = C_{C_a(b)}(x)$.
\end{proposition}

\begin{proof}
Let $A(t) = \sum_{i \geq 1} a_i t^i / i!$ and $X(t) = \sum_i x_i t^i / i!$. By
Lemma~\ref{lem:egf}, the exponential generating function of $C_a(x)$ is
$\sum_i x_i A(t)^i / i! = X(A(t))$. The partial Bell polynomials do not read the
entry of index $0$, and the entries of $C_a(b)$ from index $1$ on have generating
function $Y(A(t))$, where $Y(t) = \sum_{i \geq 1} b_i t^i / i!$. So both sides
have generating function $X(Y(A(t)))$, because composition of formal power series
without constant term is associative.
\end{proof}

The second equation for the compositional inverse, $f[g(t) - 1] = 1 + t$, also
holds \cite{DiNardoSenato2001}. Its proof needs associativity of
composition.

\begin{corollary}[The Other Side of the Inverse]
\label{cor:inverse-right}
Let $a$ be a sequence with $a_1 = 1$, and let $g$ be its inverse moments. Then
$\sum_{i \leq m} a_i\, B_{m,i}(g) = \chi_m$ for every $m \geq 1$, that is,
$f[g(t) - 1] = 1 + t$, where $f(t) = \sum_m a_m t^m / m!$ and
$g(t) = \sum_m g_m t^m / m!$.
\end{corollary}

\begin{proof}
Proposition~\ref{prop:inverse} says $C_a(g) = \chi$. Two facts about $\chi$ are
needed. First, $B_{m,i}(\chi)$ is $1$ for $i = m$ and $0$ otherwise, because the
only partition all of whose parts have size $1$ is the partition into singletons.
Hence $C_\chi(x) = x$ for every $x$. Second, $C_a(\chi)_m = B_{m,1}(a) = a_m$ for
$m \geq 1$.

The equation of Proposition~\ref{prop:inverse} at $m = 1$ gives $g_1 = 1$, so $g$
has inverse moments $h$, with $C_g(h) = \chi$. By Proposition~\ref{prop:assoc},
\[
  h = C_\chi(h) = C_{C_a(g)}(h) = C_a(C_g(h)) = C_a(\chi) ,
\]
so $h_m = a_m$ for $m \geq 1$. Let $m \geq 1$. Then $B_{m,0}(g) = 0$, so the entry
of index $0$ does not enter, and
$\sum_{i \leq m} a_i\, B_{m,i}(g) = \sum_{i \leq m} h_i\, B_{m,i}(g) = \chi_m$.
\end{proof}

\section{Naturality versus Parametricity}
\label{sec:separation}

In the previous chapter we have seen that everything collapsed to naturality. It 
is not true anymore.

\begin{definition}
\label{def:wdict}
A \emph{$\W$-dictionary} consists of $\W$-sets $\sK$, $\sU$, $\sA$ and, for
every world $w$, a dictionary at the sets $\sK(w)$, $\sU(w)$, $\sA(w)$ in the sense
of Definition~\ref{def:class} together with a map
$\mk_w : (\NN \to \sK(w)) \to (\T\sU)(w)$, such that every method commutes with
the action of morphisms. A \emph{homomorphism} $h : d \to d'$ of $\W$-dictionaries is a triple of
natural transformations $h_K$, $h_U$, $h_A$ that commutes with every method. For
$\mk$, this means $h_{\T U}(\mk\, f) = \mk'(h_K \circ f)$, where
$h_{\T U}(\Delta, s) = (\Delta, h_U(s))$. A \emph{pair} $(d, \mathrm{dot})$
consists of a $\W$-dictionary and an implementation of the dot product at a
world, and $h$ is a \emph{homomorphism of pairs}
$(d, \mathrm{dot}) \to (d', \mathrm{dot}')$ when
\[
  h_{\T A} \circ \mathrm{dot} = \mathrm{dot}' \circ (\mathrm{id}_\NN \times h_U) .
\]
Let $F$ be a class function on pairs, with a function
$F(d, \mathrm{dot}) : \NN \times \sU \to \sK$ for each pair. It is \emph{natural} if $h_K \circ F(d, \mathrm{dot}) =
F(d', \mathrm{dot}') \circ (\mathrm{id}_\NN \times h_U)$ for every homomorphism of
pairs.
\end{definition}

The moment interpretation of Definition~\ref{def:momentdict}, with each method of
the moment dictionary at $w$ taken at the identity of $w$, is a $\W$-dictionary.
Renaming along a morphism commutes with the ring operations, with
$\mathrm{const}$ and $\mathrm{inj}$, and with $\E$ by Lemma~\ref{lem:rename}. It
carries the block allocated by $\mk$ at $w$ to the block allocated at $v$.

\begin{theorem}[Separation]
\label{thm:separation}
Let $d$ be the moment dictionary over a ring $k$ of characteristic zero, let $u$ be a 
name, and let $n \ne 1$. Then there is a natural class function $F$ on pairs with
\[
  F(d, \mathrm{dot}_{\mathrm{s}})(n, u) \ne F(d, \mathrm{dot}_{\mathrm{c}})(n, u) .
\]
Consequently:
\begin{enumerate}
\item naturality along homomorphisms of pairs does not imply that a class function takes
  the same value at $\mathrm{dot}_{\mathrm{s}}$ and $\mathrm{dot}_{\mathrm{c}}$,
  so Theorem~\ref{thm:repind} is not a consequence of naturality alone;
\item $F$ is natural and not parametric, so the analogue of
  Theorem~\ref{thm:natural} fails for the full language.
\end{enumerate}
\end{theorem}

\begin{proof}
Let $F(d, \mathrm{dot})(n, u)$ be the numeral $0 + 1 + \dots + 1$ of $\sK$, built
with the methods $0$, $1$ and $+$ of $d$, with as many summands $1$ as the block
$\Delta$ in $\mathrm{dot}\, n\, u = (\Delta, s)$ has names. Let $h$ be a
homomorphism of pairs. It acts on $\T$ by $(\Delta, s) \mapsto (\Delta, h(s))$,
so the two implementations produce blocks of the same size at corresponding
arguments. Moreover $h_K$ commutes with $0$, $1$ and $+$, so it preserves
numerals. Hence $F$ is natural. At the moment dictionary
$\mathrm{dot}_{\mathrm{s}}$ allocates one name and $\mathrm{dot}_{\mathrm{c}}$
allocates $n$, and $1 \ne n$ in $k$.

For (1), $F$ is a natural class function that takes different values at the two
implementations. For (2), the proof of Theorem~\ref{thm:repind} shows that the
two implementations are related at $\Nat \to U \to \T A$. A parametric class function
would send them to values related at $K$, which is equality, and $F$ does not.
\end{proof}

\chapter{Formalization and Implementation}
\label{ch:formal}

The Lean development and the Haskell library are in the repository
\url{https://github.com/sonatsuer/umbral-calculus}. The Lean~4 development in
\texttt{FreeThms} proves the results of
Chapters~\ref{ch:restricted} and~\ref{ch:full}, using Mathlib \cite{Mathlib}. It
has no unproved statements: \texttt{scripts/check.sh} builds it and fails on any
\texttt{sorry}. The file \texttt{FreeThms/Statements.lean} restates every result
of Table~\ref{tab:correspondence}. The Haskell library in \texttt{haskell}
implements both languages and checks the results on examples, in the moment
model, symbolically, modulo $n$ and as an expectation.

Table~\ref{tab:correspondence} lists, for each result, the Lean theorem that
proves it, with its name in \texttt{Statements.lean} in brackets, and the Haskell
code and the group of checks in \texttt{test/Main.hs} that exercise it. Lean
files are relative to \texttt{FreeThms}. A dash means that there is no
counterpart.

{\footnotesize
\begin{longtable}{@{}>{\raggedright\arraybackslash}p{2.4cm}>{\raggedright\arraybackslash}p{6.1cm}>{\raggedright\arraybackslash}p{4.9cm}@{}}
\caption{Results, Lean theorems and Haskell code.}\label{tab:correspondence}\\
\toprule
Result & Lean (file: theorem [statement]) & Haskell (code; checks) \\
\midrule
\endfirsthead
\toprule
Result & Lean (file: theorem [statement]) & Haskell (code; checks) \\
\midrule
\endhead
\bottomrule
\endfoot
\multicolumn{3}{@{}l}{\emph{Chapter~\ref{ch:intro}}} \\
Def.~\ref{def:umbral} & \texttt{Umbral/Moments}: \texttt{UmbralCalculus} & -- \\
Prop.~\ref{prop:bernoulli-rec} & -- & \texttt{bernoulliDefect}; \texttt{bernoulli} \\
Prop.~\ref{prop:euler} & -- & \texttt{eulerTerm}, \texttt{pitfallTerm}; \texttt{showcaseChecks} \\
Def.~\ref{def:moment}, Prop.~\ref{prop:moment} & \texttt{Umbral/Moments}: \texttt{eval}, \texttt{moments\_eq}, \texttt{momentModel} [\texttt{moment\_model\_is\_umbral\_calculus}]; converse: \texttt{UmbralCalculus.eq\_eval} [\texttt{umbral\_calculus\_is\_moment\_model}] & \texttt{momentModel}, \texttt{evaluate} \\
Def.~\ref{def:probmodel}, Prop.~\ref{prop:probmodel} & \texttt{Umbral/Probabilistic}: \texttt{IsModel}, \texttt{realize}, \texttt{probModel}, \texttt{integral\_realize} [\texttt{probabilistic\_model\_is\_moment\_model}] & \texttt{probabilityModel}, \texttt{exact}; \texttt{probability} \\
\midrule
\multicolumn{3}{@{}l}{\emph{Chapter~\ref{ch:restricted}}} \\
Defs.~\ref{def:types}--\ref{def:logrel} & \texttt{Core/SystemF}: \texttt{PTy}, \texttt{PTm}, \texttt{PTyD}, \texttt{peval}, \texttt{PTyRel} & -- \\
Thm.~\ref{thm:abstraction-precise} & \texttt{Core/SystemF}: \texttt{pabstraction} & -- \\
Def.~\ref{def:class}, type~\eqref{eq:type} & \texttt{Umbral/Language}: \texttt{Dict}, \texttt{oneUmbraTy}, \texttt{den} & \texttt{Umbral}, \texttt{Expr1} \\
Thm.~\ref{thm:free} & \texttt{Umbral/Language}: \texttt{free\_quantTy}, \texttt{applyN\_rel}, \texttt{denN\_rel} [\texttt{umbral\_free\_theorem\_many}]; $m = 1$: \texttt{den\_rel} [\texttt{umbral\_free\_theorem}] & -- \\
Lem.~\ref{lem:rename} & \texttt{Umbral/Moments}: \texttt{eval\_rename} [\texttt{eval\_preserved\_by\_moment\_maps}] & -- \\
Thm.~\ref{thm:substitution} & \texttt{Umbral/Language}: \texttt{substitution} [\texttt{substitution\_principle\_many}]; $m = 0$: \texttt{similarity} [\texttt{substitution\_principle}] & \texttt{oneUmbraExamples}, \texttt{twoUmbraExamples}; \texttt{substitution} \\
Thm.~\ref{thm:uncorrelation} & \texttt{Umbral/Language}: \texttt{moments\_of\_common\_target}, \texttt{no\_common\_target} [\texttt{uncorrelation\_constrains\_relations}]; \texttt{Umbral/Moments}: \texttt{injectivity\_necessary} & \texttt{productMoment}; \texttt{uncorrelation} \\
Thm.~\ref{thm:exchange} & \texttt{Umbral/Language}: \texttt{den2\_rel}, \texttt{exchangeable} [\texttt{similar\_umbrae\_exchangeable}] & \texttt{twoUmbraExamples}; \texttt{exchangeability} \\
Thm.~\ref{thm:integral} & \texttt{Umbral/Language}: \texttt{respects\_subring}, \texttt{integral} [\texttt{integrality}] & -- \\
Def.~\ref{def:natural}, Thm.~\ref{thm:natural} & \texttt{Umbral/Naturality}: \texttt{Family}, \texttt{Parametric}, \texttt{Natural}, \texttt{parametric\_iff\_natural} [\texttt{parametric\_is\_natural}] & -- \\
Cor.~\ref{cor:natural} & \texttt{Umbral/Naturality}: \texttt{den\_natural} [\texttt{terms\_are\_natural}]; the readings: \texttt{similarity\_natural}, \texttt{DictHom.ofSubring}, \texttt{integrality\_natural} [\texttt{substitution\_principle\_natural}, \texttt{integrality\_as\_naturality}] & \texttt{Mod}; \texttt{modular} (along $\ZZ \to \ZZ/97$) \\
Thm.~\ref{thm:probabilistic} & \texttt{Umbral/Probabilistic}: \texttt{realizeHom}, \texttt{probabilistic\_soundness} [same name] & \texttt{probabilityModel}, \texttt{exact}, \texttt{monteCarlo}; \texttt{probability} \\
\midrule
\multicolumn{3}{@{}l}{\emph{Chapter~\ref{ch:full}}} \\
Def.~\ref{def:inverse}, Prop.~\ref{prop:inverse} & \texttt{Umbral/Inverse}: \texttt{inverseT}, \texttt{sem\_inverse}, \texttt{inverse\_spec}, \texttt{inverse\_unique} [\texttt{inverse\_umbra}, \texttt{inverse\_composition}, \texttt{inverse\_uniqueness}] & \texttt{inverse}, \texttt{inverseMoments}; \texttt{inverseChecks} \\
Defs.~\ref{def:world}--\ref{def:momentdict} & \texttt{Core/Kripke}: \texttt{AllocStr}, \texttt{QTy}, \texttt{QTyD}, \texttt{eval}; \texttt{Umbral/Worlds}: \texttt{World}, \texttt{Hom}, \texttt{alloc}, \texttt{momentMk}, \texttt{applyMethods}, \texttt{applyDict} & \texttt{Fresh}, \texttt{Name}, \texttt{momentModel} \\
Defs.~\ref{def:worldrel}, \ref{def:future} & \texttt{Umbral/Worlds}: \texttt{WRel}, \texttt{Fut} & -- \\
Lem.~\ref{lem:extension} & \texttt{Umbral/Worlds}: \texttt{extension} [\texttt{extension\_lemma}] & -- \\
Def.~\ref{def:kripke-logrel} & \texttt{Core/Kripke}: \texttt{RelStr}, \texttt{QRel} & -- \\
Lem.~\ref{lem:monotone} & \texttt{Core/Kripke}: \texttt{QRel.mono} & -- \\
Lem.~\ref{lem:dictrel} & \texttt{Umbral/Worlds}: \texttt{applyDict\_rel}, \texttt{mk\_rel} & -- \\
Thm.~\ref{thm:abstraction2} & \texttt{Core/Kripke}: \texttt{abstraction}; free theorem: \texttt{Umbral/Worlds}: \texttt{den\_rel} [\texttt{umbral\_free\_theorem\_alloc}]; substitution: \texttt{substitution} [\texttt{substitution\_principle\_full}] & \texttt{allocatingExamples}; \texttt{substitutionWithAllocation} \\
Prop.~\ref{prop:conservative} & \texttt{Core/Conservativity}: \texttt{conservative}; \texttt{Umbral/Conservativity}: \texttt{conservative} [\texttt{full\_language\_conservative}] & -- \\
Lem.~\ref{lem:dotmoments} & \texttt{Umbral/DotMoments}: \texttt{eval\_copies}, \texttt{bell\_dotMoments} [\texttt{dot\_moments\_copies}, \texttt{dot\_moments\_bell}] & \texttt{dotMoments}, \texttt{bell}; \texttt{dotMomentsCopies}, \texttt{dotMomentsBell} \\
Def.~\ref{def:auxiliary}, Prop.~\ref{prop:auxiliary}(1)--(3) & \texttt{Umbral/AuxiliaryUmbrae}: \texttt{newT}, \texttt{dotConstructedT}, \texttt{dotSymbolicT}, \texttt{sem\_new}, \texttt{sem\_dotConstructed}, \texttt{sem\_dotSymbolic} [\texttt{fresh\_copy}, \texttt{dot\_constructed}, \texttt{dot\_symbolic}] & \texttt{new}, \texttt{dotConstructed}, \texttt{dotSymbolic}; \texttt{freshCopy}, \texttt{dotConstructedChecks}, \texttt{dotSymbolicChecks} \\
Prop.~\ref{prop:auxiliary}(4) & \texttt{Umbral/Composition}: \texttt{composeT}, \texttt{sem\_compose}, \texttt{composeMoments\_eq}, \texttt{composeMoments\_const} [\texttt{composition}, \texttt{composition\_moments}, \texttt{composition\_constant}] & \texttt{compose}; \texttt{composition}, \texttt{compositionMoments}, \texttt{compositionConstant}, \texttt{compoundPoissonChecks} \\
Prop.~\ref{prop:auxiliary}(5) & \texttt{Umbral/Composition}: \texttt{singletonT}, \texttt{cumulantT}, \texttt{sem\_cumulant} [\texttt{cumulant\_umbra}] & \texttt{singleton}, \texttt{cumulant}; \texttt{cumulantUmbra} \\
Def.~\ref{def:similar-comp} & \texttt{Umbral/Similarity}: \texttt{resultMoments}, \texttt{resultMomentsA} & -- \\
Prop.~\ref{prop:similar} & \texttt{Umbral/Similarity}: \texttt{new\_similar}, \texttt{dotSymbolic\_similar}, \texttt{dotConstructed\_similar}, \texttt{cumulant\_similar}, \texttt{inverse\_similar} [\texttt{auxiliary\_umbrae\_similar}]; \texttt{compose\_similar}, \texttt{composeA\_similar} [\texttt{composition\_similar}, \texttt{composition\_polynomial\_similar}] & -- \\
Def.~\ref{def:interface} & \texttt{Umbral/RepIndependence}: \texttt{CtxTerm}, \texttt{denCtx}; \texttt{Umbral/AuxiliaryUmbrae}: \texttt{Prog} & \texttt{CtxDot}, \texttt{CtxNew}, \texttt{CtxSum} \\
Thm.~\ref{thm:repind} & \texttt{Umbral/RepIndependence}: \texttt{repind}, \texttt{repind\_dot} [\texttt{representation\_independence}, \texttt{repind\_dot\_product}] & \texttt{dotContexts}; \texttt{dotProductIndependence} \\
Cor.~\ref{cor:repind-new} & \texttt{Umbral/RepIndependence}: \texttt{repind\_new} [\texttt{repind\_fresh\_copy}] & \texttt{newAfterUnused}, \texttt{newContexts}; \texttt{freshCopyIndependence} \\
Prop.~\ref{prop:distributive} & \texttt{Umbral/Distributivity}: \texttt{compTransform\_conv} [\texttt{distributivity}] & \texttt{composeA}; \texttt{distributivity} \\
Cor.~\ref{cor:repind-compose} & \texttt{Umbral/RepIndependence}: \texttt{repind\_composeSum} [\texttt{repind\_composition\_of\_sum}] & \texttt{composeSumFused}, \texttt{composeSumSplit}, \texttt{sumContexts}; \texttt{compositionOfSumIndependence} \\
Cor.~\ref{cor:sheffer} & \texttt{Umbral/RepIndependence}: \texttt{sheffer} [\texttt{sheffer\_identity}] & --; \texttt{shefferUmbral} \\
Prop.~\ref{prop:assoc} & \texttt{Umbral/Distributivity}: \texttt{compTransform\_succ}, \texttt{pBell\_compTransform}, \texttt{compTransform\_assoc} [\texttt{composition\_associative}] & -- \\
Cor.~\ref{cor:inverse-right} & \texttt{Umbral/Inverse}: \texttt{inverse\_spec\_right} [\texttt{inverse\_composition\_right}] & \texttt{inverse}; \texttt{inverseChecks} \\
Def.~\ref{def:wdict} & \texttt{Umbral/Separation}: \texttt{WSet}, \texttt{WSet.T}, \texttt{WDict}, \texttt{WDictHom}, \texttt{DotPair}, \texttt{DotPairHom}, \texttt{PairFamily.Natural}, \texttt{momentWDict} & -- \\
Thm.~\ref{thm:separation} & \texttt{Umbral/Separation}: \texttt{blockCount}, \texttt{blockCount\_natural}, \texttt{separation} [\texttt{separation\_from\_naturality}] & \texttt{countFresh}; \texttt{allocationCount} \\
\end{longtable}
}

\begin{thebibliography}{99}

\bibitem{Aldous1985}
D.~J. Aldous,
\emph{Exchangeability and related topics},
in: \'Ecole d'\'Et\'e de Probabilit\'es de Saint-Flour XIII -- 1983,
Lecture Notes in Mathematics 1117, Springer, 1985, pp.~1--198.

\bibitem{Blissard1861}
J.~Blissard,
\emph{Theory of generic equations},
Quart. J. Pure Appl. Math. \textbf{4} (1861), 279--305.

\bibitem{DiNardoNiederhausenSenato}
E.~Di~Nardo, H.~Niederhausen and D.~Senato,
\emph{The classical umbral calculus: Sheffer sequences},
arXiv:0810.3554.

\bibitem{DiNardoSenato2001}
E.~Di~Nardo and D.~Senato,
\emph{Umbral nature of the Poisson random variables},
in: Algebraic Combinatorics and Computer Science, Springer, 2001, pp.~245--266;
arXiv:math/0412054.

\bibitem{DiNardoSenato2006}
E.~Di~Nardo and D.~Senato,
\emph{An umbral setting for cumulants and factorial moments},
European J. Combin. \textbf{27} (2006), 394--413; arXiv:math/0412052.

\bibitem{DiNardoSenatoKstat}
E.~Di~Nardo and D.~Senato,
\emph{A symbolic method for $k$-statistics},
arXiv:math/0505033.

\bibitem{Gessel2003}
I.~M. Gessel,
\emph{Applications of the classical umbral calculus},
Algebra Universalis \textbf{49} (2003), 397--434.

\bibitem{Mathlib}
The mathlib Community,
\emph{The Lean mathematical library},
in: Proceedings of the 9th ACM SIGPLAN International Conference on Certified
Programs and Proofs (CPP 2020), 2020, pp.~367--381.

\bibitem{OHearnTennent1995}
P.~W. O'Hearn and R.~D. Tennent,
\emph{Parametricity and local variables},
J. ACM \textbf{42} (1995), 658--709, doi:10.1145/210346.210425.

\bibitem{PittsStark1993}
A.~M. Pitts and I.~D.~B. Stark,
\emph{Observable properties of higher order functions that dynamically create
local names, or: What's new?},
in: Mathematical Foundations of Computer Science (MFCS 1993), Lecture Notes in
Comput. Sci. 711, Springer, 1993, pp.~122--141.

\bibitem{Reynolds1983}
J.~C. Reynolds,
\emph{Types, abstraction and parametric polymorphism},
in: Information Processing 83, North-Holland, 1983, pp.~513--523.

\bibitem{Reynolds1984}
J.~C. Reynolds,
\emph{Polymorphism is not set-theoretic},
in: Semantics of Data Types, Lecture Notes in Comput. Sci. 173, Springer, 1984,
pp.~145--156.

\bibitem{RotaTaylor1994}
G.-C. Rota and B.~D. Taylor,
\emph{The classical umbral calculus},
SIAM J. Math. Anal. \textbf{25} (1994), 694--711.

\bibitem{Wadler1989}
P.~Wadler,
\emph{Theorems for free!},
in: Proceedings of the Fourth International Conference on Functional Programming
Languages and Computer Architecture (FPCA '89), 1989, pp.~347--359.

\end{thebibliography}

\end{document}
